

\documentclass[11pt]{article}

%--- Packages (kept minimal; mirrors BCF.tex where portable) ------------
\usepackage[margin=1in]{geometry}
\usepackage{amsmath, amssymb, amsthm, amsfonts, bm}
\usepackage{booktabs, multirow}
\usepackage{natbib}
\usepackage{algorithm, algpseudocode}
\usepackage{enumitem}
\usepackage{graphicx}
\usepackage{subcaption}
\usepackage[colorlinks, citecolor=blue, urlcolor=blue, linkcolor=blue]{hyperref}
\graphicspath{{../tests/analysis/plots/}}

%--- Theorem environments ----------------------------------------------
\theoremstyle{plain}
\newtheorem{thm}{Theorem}[section]
\newtheorem{prop}[thm]{Proposition}
\newtheorem{lemma}[thm]{Lemma}
\theoremstyle{definition}
\newtheorem{defn}[thm]{Definition}
\newtheorem{assumption}[thm]{Assumption}
\theoremstyle{remark}
\newtheorem{rem}[thm]{Remark}

\newcommand\independent{\protect\mathpalette{\protect\independenT}{\perp}}
\def\independenT#1#2{\mathrel{\rlap{$#1#2$}\mkern2mu{#1#2}}}

%--- Distributions / shorthand -----------------------------------------
\newcommand{\iid}{\stackrel{\mbox{\tiny iid} }{\sim}}
\newcommand{\logit}{\mathrm{logit}}
\newcommand{\Poi}{\mathrm{Poisson}}
\newcommand{\NB}{\mathrm{NB}}
\newcommand{\Ber}{\mathrm{Bernoulli}}
\newcommand{\Be}{\mathrm{Beta}}
\newcommand{\Ga}{\mathrm{Gamma}}
\newcommand{\GIG}{\mathrm{GIG}}
\newcommand{\N}{\mathrm{N}}
\newcommand{\Exp}{\mathrm{Exponential}}

%--- Vectors / matrices --------------------------------------------------
\newcommand{\X}{\mathbf{X}}
\newcommand{\Y}{\mathbf{Y}}
\newcommand{\Z}{\mathbf{Z}}
\newcommand{\x}{\mathrm{x}}
\newcommand{\z}{\mathrm{z}}
\newcommand{\y}{\mathrm{y}}
\newcommand{\I}{\mathbf{I}}
\newcommand{\bbeta}{\boldsymbol{\beta}}
\newcommand{\btheta}{\boldsymbol{\theta}}
\newcommand{\bSigma}{\boldsymbol{\Sigma}}

%--- Probability / expectation operators --------------------------------
\newcommand{\E}{\mathbb{E}}
\newcommand{\Var}{\mathrm{Var}}
\newcommand{\Cov}{\mathrm{Cov}}
\renewcommand{\Pr}{\mathrm{Pr}}
\newcommand{\ind}{\mathbf{1}}

%--- BCF-style (mu, tau) decomposition: count component -----------------
\newcommand{\muF}{\mu_f}
\newcommand{\tauF}{\tau_f}
\newcommand{\muz}{\mu_{z}}
\newcommand{\muone}{\mu_{1}}
\newcommand{\muzero}{\mu_{0}}
\newcommand{\tauresp}{\tau}                 % response-scale CATE
\newcommand{\taulog}{\tau_{f}}              % log-rate-scale CATE (= tauF)
\newcommand{\pihat}{\widehat{\pi}}

%--- GIG leaf-prior reparameterisation ----------------------------------
\newcommand{\azero}{a_{0}}
\newcommand{\azerotau}{a_{0,\tau}}
\newcommand{\leafc}{c}
\newcommand{\leafd}{d}

%--- ZI-specific notation -----------------------------------------------
\newcommand{\Si}{S_{i}}
\newcommand{\pzi}[1][z]{p_{\mathrm{zi},\,#1}}
\newcommand{\etazi}{\eta_{\mathrm{zi}}}
\newcommand{\Detazi}{\Delta\etazi}
\newcommand{\fzeroSym}{f_{0}}
\newcommand{\foneSym}{f_{1}}
\newcommand{\muFzero}{\mu_{f_{0}}}
\newcommand{\muFone}{\mu_{f_{1}}}
\newcommand{\tauFzero}{\tau_{f_{0}}}
\newcommand{\tauFone}{\tau_{f_{1}}}
\newcommand{\zconc}{z_{\mathrm{conc}}}
\newcommand{\zconctau}{z_{\mathrm{conc},\,\tau}}
\newcommand{\zilat}{z_{i}}
\newcommand{\xilat}{\xi_{i}}
\newcommand{\philat}{\phi_{i}}
\newcommand{\omegai}{\omega_{i}}
\newcommand{\logfgroup}[1]{\log f^{(#1)}}
\newcommand{\logsumexp}{\mathrm{logsumexp}}
\newcommand{\softw}{w}
\newcommand{\Deltalam}{\Delta\lambda}
\newcommand{\Deltapzi}{\Delta p_{\mathrm{zi}}}

% =======================================================================
\title{CountBCF: Bayesian Causal Forests for Count and
       Zero-Inflated Count Outcomes}
\author{Hugo Gobato Souto\thanks{Dell Technologies. ORCID: \href{https://orcid.org/0000-0002-7039-0572}{0000-0002-7039-0572}.
Email: \href{mailto:hugo.souto@dell.com}{hugo.souto@dell.com}.}}
\date{May 2026}

\begin{document}
\maketitle

\begin{abstract}
\noindent
Count outcomes -- hospital admissions, infection cases, insurance
claims, citation counts, recidivism events, single-cell read counts --
are pervasive in the applied questions for which heterogeneous
treatment effect (HTE) methods were designed, and a non-trivial
fraction of these outcomes are also \emph{zero-inflated}: some units
never enter the count-generating process at all (a structural zero),
while the rest accumulate strictly positive counts at heterogeneous
intensities. A treatment can move the structural-zero probability and
the count intensity independently, and the response-scale conditional
average treatment effect (CATE) mixes the two channels. Existing
tools collapse this structure: the workhorse Bayesian Causal Forest
(BCF) of \citet{hahn2020bayesian} assumes a Gaussian response;
Poisson and negative binomial generalised linear models, doubly robust
extensions, and generalised random forests either rule out per-unit
CATEs by design or discard BCF's two structural ingredients (the
differentially shrunk $(\mu, \tau)$ decomposition and the
propensity-as-covariate prognostic forest that combats
regularization-induced confounding); classical zero-inflated
regressions assume rigid treatment interactions; two-part learners
separate the channels but discard the shared-covariate information.
This paper introduces \textbf{CountBCF}, which embeds the BCF
decomposition inside the log-linear BART backbone of
\citet{murray2021loglinear}. The framework covers two regimes within a
single model: a \emph{count} regime that delivers a $(\muF, \tauF)$
pair on the log-rate scale, and a \emph{zero-inflated} regime that
extends the decomposition to all three log-linear-BART forests --
$f, \fzeroSym, \foneSym$ -- yielding six forests in total. Each
function group inherits the BCF differential-shrinkage philosophy and
the propensity-as-covariate fix; posterior draws translate to per-unit
response-scale CATEs in closed form, with a per-unit channel
decomposition into count-rate and structural-zero effects available
in the zero-inflated case. We release the \texttt{countbcf} R package,
built on the \texttt{countbart} MCMC backend of
\citet{wikle2023causal}, and pre-register a simulation study that
crosses linear and nonlinear data-generating processes with Poisson,
negative binomial, ZIP, and ZINB likelihoods and sweeps one-at-a-time
over sample size, covariate dimension, and target average treatment
effect.
\end{abstract}

\noindent\textbf{Keywords:} Bayesian causal forests; Conditional average
treatment effects; Count regression; Heterogeneous treatment effects;
Log-linear BART; Poisson and negative binomial regression;
Zero-inflated regression; Structural zeros.

% =======================================================================
\section{Introduction}\label{sec:intro}

Many of the outcomes that motivate modern causal inference are
counts. Hospital admissions and emergency-department visits,
infection incidences, dental caries, traffic crashes, insurance
claims, citation totals, recidivism events, household child counts,
and single-cell unique-molecular-identifier reads are all
integer-valued and frequently bunched at zero. The applied
literatures that produced these outcomes -- epidemiology, health
economics, ecology, operations research, finance, social statistics,
computational biology -- share an interest not just in whether an
intervention works on average, but in \emph{for whom} it works, and
by how much: the conditional average treatment effect, or CATE. The
past decade of heterogeneous treatment effect (HTE) methodology has
answered that demand with flexible nonparametric tools -- Bayesian
Additive Regression Trees \citep{chipman2010bart, hill2011bayesian},
Bayesian Causal Forests \citep[BCF;][]{hahn2020bayesian}, generalised
random forests \citep{athey2019generalized}, meta-learners
\citep{kunzel2019metalearners}, and double machine learning
\citep{chernozhukov2018double} -- but the bulk of that methodology
assumes a continuous, approximately Gaussian response.

Treating counts as if they were Gaussian has three well-documented
failure modes \citep{camerontrivedi2013count, murray2021loglinear}.
Count distributions place positive mass at zero and at small
integers, often on a large fraction of the support; the conditional
variance grows with the conditional mean, often super-linearly under
overdispersion; and the substantive scientific quantity is almost
always a rate or a relative risk, both of which live on the log
scale. These pathologies bias point estimates, distort uncertainty
quantification, and degrade the heterogeneity signal that HTE methods
are designed to recover. A second, sharper failure occurs when the
count outcome is also \emph{zero-inflated}: a non-trivial fraction of
units never enter the count-generating process at all (a structural
zero), while the rest accumulate strictly positive counts at
heterogeneous intensities. A canonical example is annual healthcare
utilisation: in many provider panels and insurance cohorts, a
sizeable fraction of enrollees record zero visits in a given year
because they are effectively out of the care-seeking system, while
the remaining enrollees record strictly positive visit counts at
heterogeneous intensities. The same pattern arises in manufacturing
defect counts \citep{lambert1992zip}, automobile insurance claims,
ecological species-abundance surveys, sentencing and recidivism
counts, and single-cell RNA sequencing, where a large excess-zero
mass reflects technical dropout on top of a count process.

What makes zero-inflation especially awkward for causal inference is
that an intervention can move \emph{both} mechanisms, and the two
movements need not point in the same direction. Continuing the
healthcare example: a workplace wellness programme aimed at preventive
care could plausibly pull some enrollees \emph{into} the system --
shrinking the structural-zero fraction -- while simultaneously
\emph{lowering} the visit intensity of those already engaged because
preventive care substitutes for acute episodes. The response-scale
CATE is then a mixture of the two channel effects; whether it ends
up positive, negative, or near zero depends on which channel
dominates for which unit. A method that ignores the structural-zero
mechanism necessarily mis-attributes its movement: if the treatment
shifts the structural-zero probability but leaves the count rate
untouched, a Poisson-only working model has to absorb the entire
change into the rate estimate, biasing it; the same mis-attribution
flows in the opposite direction when the treatment moves the count
rate but not the structural-zero probability.

Faced with a count outcome and a CATE question, an applied analyst
today reaches for one of three default tools, each of which gives up
something the analysis arguably needs. The first is a Gaussian HTE
method, most commonly BCF or causal BART, applied to the raw counts
as if they were continuous; this preserves the regularisation and
uncertainty quantification practitioners trust, but commits all three
Gaussian pathologies above. The second is a parametric Poisson,
negative binomial, ZIP, or ZINB generalised linear model with
treatment-by-covariate interactions \citep{camerontrivedi2013count,
lambert1992zip, greene1994zinb}, possibly stabilised by propensity
weighting \citep{rosenbaum1983central} or a doubly robust correction
\citep{chernozhukov2018double}; this respects the count likelihood
but recovers heterogeneity only through pre-specified interaction
terms, defeating the nonparametric premise that motivated the move
to HTE methods in the first place. The third is a generalised random
forest with a Poisson moment condition \citep{athey2019generalized},
or a related relative-risk-targeted estimator
\citep{kennedy2024relative}; this is honest about the count structure
but discards the two structural priors that distinguish BCF from
other forest-based HTE methods, leaving heterogeneity estimates
exposed to the high-variance behaviour that motivated those priors
to begin with. Two-part and hurdle learners \citep{mullahy1986hurdle}
separate the channels of a zero-inflated outcome but do not share
covariate information across them and do not implement BCF
differential shrinkage.

The structural priors in question are the distinguishing features of
\citet{hahn2020bayesian}. Their construction decomposes the
conditional mean of the response into a \emph{prognostic} function --
what the outcome would be in the absence of treatment, as a smooth
function of the covariates -- and a \emph{moderating} function -- by
how much treatment shifts that baseline at each point in covariate
space. The two functions are fit by separate BART ensembles with
deliberately different priors: the prognostic forest is large and
flexible; the moderating forest is small and tightly shrunk toward
zero, reflecting the prior belief that treatment effect heterogeneity
is typically smoother and weaker than the prognostic surface itself.
This differential shrinkage stabilises CATE estimates in finite
samples without forcing them to be constant. The second structural
ingredient is the inclusion of an estimated propensity score among
the inputs to the prognostic forest only. \citet{hahn2020bayesian}
show that omitting this control creates \emph{regularization-induced
confounding} (RIC): the moderating forest absorbs prognostic signal
that the regularised prognostic forest underfits, biasing CATE
estimates in the direction of the propensity score. Both pathologies
are at least as severe under a log-link count likelihood as under a
Gaussian one: the multiplicative geometry of the log scale amplifies
the consequences of prognostic underfitting, and the count
likelihood's tight mean-variance coupling leaves the moderating
forest doubly exposed to RIC bias. The same arguments apply
\emph{per channel} in the zero-inflated case.

\medskip
\noindent\textbf{Our contribution.}
This paper develops \textbf{CountBCF}, a unified Bayesian causal
forest framework for count and zero-inflated count outcomes that
retains both of BCF's structural ingredients while replacing its
Gaussian backbone with a log-linear count likelihood. The framework
covers two regimes within a single model and a single sampler:
\begin{enumerate}
  \item \textbf{A count-faithful BCF.}
    For non-zero-inflated count outcomes, we extend the Bayesian
    Causal Forest of \citet{hahn2020bayesian} from a Gaussian to a
    Poisson and a negative binomial likelihood by routing both the
    prognostic and the moderating forest through the log-linear BART
    backbone of \citet{murray2021loglinear}. The result is a BCF
    that respects discreteness, mass at zero, and mean-variance
    coupling without abandoning the prognostic/moderating
    decomposition that defines the BCF family.
  \item \textbf{A six-forest BCF decomposition for zero-inflated
    outcomes.}
    Murray's log-linear BART parameterises a zero-inflated count
    likelihood through three function groups: the count log-rate
    $f$, the structural-zero log-odds $\fzeroSym$, and the
    not-structural-zero log-odds $\foneSym$, with the structural-zero
    probability identified as $\sigma(\fzeroSym - \foneSym)$. We
    replace each group's single forest with the BCF
    propensity-adjusted prognostic / moderator pair, yielding six
    forests in total: $(\muF, \tauF)$ for the count channel and
    $(\muFzero, \tauFzero), (\muFone, \tauFone)$ for the ZI channel.
    The non-zero-inflated case is recovered as the three-forest sub-model.
  \item \textbf{RIC protection on the log scale, per channel.}
    We retain the propensity-as-covariate fix and the
    differential-shrinkage prior structure of BCF, with the
    prognostic forest of each function group receiving more trees and
    a weaker zero-bias prior than the corresponding moderating
    forest. The leaf prior is the GIG mixture of
    \citet{murray2021loglinear}, conjugate under a latent-variable
    augmentation; concentration hyperparameters are calibrated
    separately for the count and ZI groups, so each moderating forest
    is shrunk strictly more aggressively than its prognostic
    counterpart.
  \item \textbf{Closed-form response-scale heterogeneity and channel
    decomposition.}
    The log-linear parameterisation means that, given posterior draws
    of the forests at any covariate value, the response-scale
    conditional treatment effect is available without any additional
    Monte Carlo step: it is a deterministic function of the stored
    forest values. We report per-unit posterior means, $50\%$ and
    $95\%$ credible intervals, and population-averaged ATE summaries,
    all on the response scale. In the zero-inflated case the same
    posterior draws produce a per-unit decomposition of the CATE into
    a count-rate channel $\Deltalam(x)$ and a structural-zero channel
    $\Deltapzi(x)$.
  \item \textbf{Open software and a pre-registered simulation study.}
    We release \texttt{countbcf}, an open-source R package built on
    the \texttt{countbart} MCMC backend of \citet{wikle2023causal},
    with a single user-facing entry point that consumes a
    treated/control split and returns the posterior draw matrices
    needed for the closed-form CATE. The package is accompanied by a
    pre-registered simulation study (Section~\ref{sec:methodology})
    that crosses linear and nonlinear data-generating processes with
    Poisson, NB, ZIP, and ZINB likelihoods and sweeps one-at-a-time
    over sample size, covariate dimension, and target average
    treatment effect, scoring all methods against a rate-scale truth.
\end{enumerate}

The remainder of the paper proceeds as follows.
Section~\ref{sec:litreview} places CountBCF in the literature on
count and zero-inflated causal inference and identifies the specific
gap it fills. Section~\ref{sec:notation} fixes notation, including
the latent structural-zero indicator and the response-scale
potential-outcome means that are the paper's target estimand.
Section~\ref{sec:background} reviews the log-linear BART machinery,
the GIG-mixture leaf prior, and the redundant $(\fzeroSym, \foneSym)$
parameterisation that makes the ZI leaf updates conditionally
conjugate. Section~\ref{sec:countbcf} introduces CountBCF proper: the
likelihood and parameterisation in both the count and the
zero-inflated regimes, the differential-shrinkage prior structure,
the augmented Gibbs sampler, and the closed-form response-scale CATE
and channel-decomposition summaries.
Section~\ref{sec:methodology} pre-registers the simulation study used
to evaluate CountBCF against a Gaussian-BCF baseline.
Sections~\ref{sec:empirical} and \ref{sec:discussion} are reserved
for empirical results and discussion, to be written once the
simulation study completes.

% =======================================================================
\section{Relationship to previous literature}\label{sec:litreview}

This section reviews the literatures that bear on conditional
average treatment effect (CATE) estimation for count and
zero-inflated count outcomes. The cell of the methodological grid in
which CountBCF sits -- causal $\times$ count-aware (with or without
zero-inflation) $\times$ nonparametric Bayesian -- is sparsely
populated; the subsections that follow walk the adjacent cells in
turn before Section~\ref{sec:lit-gap} states the gap that CountBCF
fills.

\subsection{Causal inference for count outcomes:
parametric and semiparametric eras}\label{sec:lit-count}

Causal analysis of count outcomes has its longest tradition in the
parametric era of generalised linear models. The Poisson and negative
binomial GLMs with a log link \citep{camerontrivedi2013count} deliver
interpretable rate-ratio summaries of the average causal effect under
an unconfoundedness assumption and a correctly specified mean
function. These models extend cleanly to the negative binomial when
equi-dispersion fails, and when unconfoundedness is plausible only
conditionally on covariates, the prescription is essentially the same
as for continuous outcomes: condition on the covariates or on the
propensity score \citep{rosenbaum1983central}. The two-stage residual
inclusion of \citet{terza2008twostage} and the marginal structural
models of \citet{robins2000marginal} relax the unconfoundedness
assumption in different ways while keeping the parametric
log-linear-mean structure. Their common feature is that they target
the \emph{average} effect under a \emph{fixed functional form}:
heterogeneity, if permitted at all, enters through pre-specified
treatment-by-covariate interactions, which is workable for a small
number of binary moderators but not for a moderate-dimensional
covariate vector with unknown interaction structure.

The semiparametric and doubly robust literature relaxes the
parametric programme's commitment to a single correct outcome model.
\citet{funk2011doubly} survey doubly robust estimators across
continuous and discrete outcomes, including Poisson outcome models
paired with logistic propensity models. \citet{van2010targetedI} and
\citet{vanderlaan2011targeted} develop targeted maximum likelihood
estimation (TMLE) as a general semiparametric framework whose
construction extends naturally to count outcomes via a Poisson or
negative binomial submodel for the targeted update step.
\citet{chernozhukov2018double} introduce the double / debiased machine
learning (DML) framework: the outcome regression and propensity
score are fit by arbitrary high-quality machine learners, cross-fitted
across folds to avoid overfitting bias, and combined in a
Neyman-orthogonal score whose root yields an asymptotically normal
ATE estimate. DML is outcome-model-agnostic: a Poisson-loss gradient
boosting machine, a neural network, or BART can sit in the outcome
slot without changing the orthogonality property. By design these
methods target the average treatment effect rather than the per-unit
CATE; heterogeneity, when reported at all, is typically delivered by
stratifying across pre-specified subgroups or by a separate post-hoc
fit of the outcome regression -- neither of which preserves the
orthogonality property that motivated the doubly robust construction.
We adopt the spirit of the doubly robust programme -- propensity
adjustment of the outcome regression -- via the route of
\citet{hansen2008prognostic} and \citet{hahn2020bayesian}, but reject
the restriction to the average effect.

\subsection{Classical and causal zero-inflated count
models}\label{sec:lit-zi}

The starting point for any zero-inflated analysis is the
zero-inflated Poisson (ZIP) regression of \citet{lambert1992zip},
which writes the response as a two-component mixture: with
probability $\pi(x)$ the unit is a structural zero, and otherwise it
is drawn from a Poisson with rate $\lambda(x)$. The hurdle
alternative \citep{mullahy1986hurdle} factors the joint density into
a binary ``any vs.\ none'' model and a zero-truncated count model;
hurdle treats \emph{all} zeros as structural while ZIP allows
sampling zeros from the count component as well.
\citet{greene1994zinb} extends both constructions to the negative
binomial case. \citet{hall2000zip} develops a random-effects ZIP for
repeated-measures data, and \citet{li2012zipidentifiability} settles
the identifiability question: in the absence of constraints, the
$(\pi, \lambda)$ pair is not jointly identified from the marginal
distribution of $Y$, but a logistic link on $\pi$ together with a log
link on $\lambda$ and at least one continuous covariate restore
identifiability. The \texttt{pscl} package of
\citet{zeileis2008regression} consolidates the count GLM toolkit --
Poisson, NB, ZIP / ZINB, hurdle -- in a single \texttt{R} treatment.

When interest shifts from describing $Y$ to estimating a treatment
effect, the two-component structure of ZI data complicates matters in
two ways. First, the response-scale estimand is itself a sum of two
interpretable contrasts: a shift in the structural-zero probability
and a shift in the count rate. Second, identifying each contrast
separately requires that the ignorability assumption be combined with
the link-function restrictions of \citet{li2012zipidentifiability} or
with an instrument acting on only one channel. The naive plug-in
strategy -- estimate a ZIP / ZINB regression with treatment-covariate
interactions and report the implied mean contrast -- inherits all of
the distributional rigidity of the underlying GLM. The two-part
learner -- fit a binary CATE for $\ind\{Y_i > 0\}$ and a count CATE
for $Y_i \mid Y_i > 0$, then combine -- correctly separates the
channels but loses joint efficiency across them and does not share
covariate information. \citet{vanderlaan2011targeted} discuss how
doubly robust / TMLE constructions adapt to these working models;
\citet{kennedy2024relative} formalises the closely related mismatch
between absolute-risk partitioning criteria and relative-risk targets.
\citet{guo2016mediation} develop an IV mediation analysis for ZI
counts that defines indirect effect \emph{ratios} suitable for
count outcomes. \citet{yu2022msnmm} propose a Multiplicative
Structural Nested Mean Model for ZI nonnegative longitudinal
outcomes. \citet{oganisian2021bayesian} build a Bayesian nonparametric
model using Dirichlet process mixtures that jointly predict
structural zeros, individualised treatment effects, and latent
subgroups. \citet{choi2023zigdag} approach ZI outcomes from the
causal-discovery direction. \citet{lukusa2017review} review the
closely related literature on zero-inflated models with missing data.
Each of these methods commits either to a particular longitudinal
structural form, to a single response-scale estimand, or to a
likelihood that does not isolate the structural-zero and count
components for inference.

\subsection{Heterogeneous-effects methods and the BART
line}\label{sec:lit-hte-bart}

Methodological work on heterogeneous treatment effects has been
dominated by Gaussian or binary outcomes. \citet{hill2011bayesian}
introduced BART \citep{chipman2010bart, chipman1998bayesian} into the
causal inference literature as a flexible learner for the outcome
regression under unconfoundedness; her construction estimates a
single BART model on $(X, Z)$ and reads the CATE off as the difference
of the model's predictions at $Z = 1$ and $Z = 0$. This estimator
inherits BART's strong empirical performance and produces a fully
Bayesian posterior over per-unit treatment effects, but it ties the
prognostic and moderating components of the conditional mean to a
single tree-structure prior, making no provision for the asymmetric
regularisation that distinguishes heterogeneity estimation from
outcome estimation, and -- applied to a count outcome -- it inherits
the three Gaussian failure modes catalogued above.

\citet{athey2019generalized} introduce generalized random forests
(GRF), a non-Bayesian alternative built on random forests with locally
weighted moment conditions. GRF accommodates a Poisson moment for
count outcomes; the asymptotic theory of \citet{wager2018estimation}
establishes pointwise consistency and asymptotic normality. GRF with
a Poisson moment shares a feature CountBCF aims to escape: the
splitting criterion targets local likelihood improvement, not
relative-risk heterogeneity. \citet{kennedy2024relative} formalises
this observation. \citet{kunzel2019metalearners} catalogue the
standard meta-learners (S, T, X, DR), which are estimator-agnostic
recipes that delegate the regression and propensity work to a
user-chosen base. In principle a meta-learner can be deployed with a
count-aware base, but the published meta-learner studies use Gaussian
or binary base learners; meta-learners do not by themselves address
RIC, and they offer no off-the-shelf solution for ZI counts.

The fourth strand -- and the one from which CountBCF draws its
computational backbone -- is the line of work on BART and its
count-aware extensions. \citet{murray2021loglinear} extends BART to
count, binomial, and multinomial outcomes by replacing the Gaussian
likelihood with a log-linear count likelihood (Poisson or negative
binomial) and replacing the Gaussian leaf prior with a generalised
inverse Gaussian (GIG) mixture leaf prior whose hyperparameters are
calibrated so that, after exponentiation, the implied prior on the
forest sum has zero mean and a user-controlled marginal variance.
The mixture structure is engineered so that, after a Gamma
latent-variable augmentation, the GIG mixture is conjugate. Section
4.2.2 of \citet{murray2021loglinear} extends the construction to
zero-inflation through the \emph{redundant} parameterisation
$\pzi(x) = \sigma\bigl(\fzeroSym(x) - \foneSym(x)\bigr)$, in which each
of $\fzeroSym$ and $\foneSym$ carries its own log-linear BART prior,
and the data drive symmetry-breaking through latent indicators. The
redundancy is precisely what makes the leaf updates conditionally
conjugate under a symmetric GIG-mixture prior. \citet{wikle2023causal}
package both constructions into the \texttt{countbart} R/C++
implementation that the present paper extends.

\citet{hahn2020bayesian} introduce Bayesian Causal Forests as the
BART-based answer to the CATE-estimation problem under
unconfoundedness, with the differential-shrinkage and
propensity-as-covariate ingredients catalogued in
Section~\ref{sec:intro}. The construction has antecedents in the
prognostic score literature of \citet{hansen2008prognostic} and the
RIC argument of \citet{hahn2016regularization}.
\citet{wikle2023causal} restrict their count BART implementation to
the predictive setting; they do not pursue the BCF decomposition.
The combination of the BCF decomposition with the log-linear count
likelihood (with or without zero-inflation), which is the content of
the present paper, has not appeared in the literature.

\subsection{Identified gap}\label{sec:lit-gap}

To summarise: the predictive count and ZI-count cells are well
populated, anchored by \citet{murray2021loglinear} and the
\texttt{countbart} implementation of \citet{wikle2023causal}. The
causal Gaussian cell is well populated, anchored by
\citet{hahn2020bayesian} on the Bayesian side and by generalized
random forests \citep{athey2019generalized, wager2018estimation} on
the frequentist side. The causal count and causal ZI-count cells are
sparsely populated, and within them no existing method combines (i)
the propensity-adjusted $(\muF, \tauF)$ decomposition that has made
BCF the de-facto Bayesian benchmark for heterogeneous causal effects,
with (ii) a count-aware (and where applicable, zero-inflation-aware)
likelihood that respects the discreteness, overdispersion, and
structural-zero process of integer outcomes. The closest prior work
either retains the Gaussian likelihood (BCF;
\citealp{hahn2020bayesian}), drops the BCF decomposition (log-linear
BART, with or without ZI; \citealp{murray2021loglinear}), separates
the two ZI channels and loses joint efficiency (two-part / hurdle
learners; \citealp{mullahy1986hurdle}; \citealp{guo2016mediation};
\citealp{yu2022msnmm}; \citealp{oganisian2021bayesian}), or replaces
the Bayesian machinery with frequentist random forests targeting an
absolute-risk or relative-risk criterion
\citep{athey2019generalized, kennedy2024relative}.
\textbf{CountBCF fills this gap.} It routes the prognostic and
moderating forest of each log-linear BART function group through the
backbone of \citet{murray2021loglinear} while preserving
\citeauthor{hahn2020bayesian}'s propensity-adjusted
differential-shrinkage prior structure, yielding a count-aware
Bayesian causal forest with closed-form per-unit response-scale
CATEs, a per-unit channel decomposition in the zero-inflated case,
and full posterior uncertainty.

% =======================================================================
\section{Problem statement and notation}\label{sec:notation}

Let $Y_i \in \mathbb{Z}_{\ge 0}$ denote a count-valued response,
$Z_i \in \{0, 1\}$ a binary treatment indicator, and
$X_i \in \mathbb{R}^p$ a length-$p$ vector of pre-treatment covariates.
We observe an iid sample $\{(Y_i, Z_i, X_i)\}_{i=1}^n$ from some joint
distribution. Capital Roman letters denote random variables, with
realisations in lower case. When $Y$ or $Z$ appear without a subscript
they denote length-$n$ column vectors; $\X$ is the $n \times p$ design
matrix of covariates.

In the potential outcomes framework \citep{imbens2015causal}, write
$Y_i(0), Y_i(1) \in \mathbb{Z}_{\ge 0}$ for the counterfactual responses
under control and treatment, with observed outcome
$Y_i = Z_i Y_i(1) + (1 - Z_i) Y_i(0)$ (SUTVA). Throughout we assume
\emph{strong ignorability}:
\begin{align}
  Y_i(0), Y_i(1) &\independent Z_i \,\big|\, X_i, \label{eq:ignorability}\\
  0 \;<\; \pi(X_i) \;&\equiv\; \Pr(Z_i = 1 \mid X_i) \;<\; 1, \label{eq:overlap}
\end{align}
the second of which (overlap) is required to identify treatment effects
everywhere in covariate space. The estimated propensity score is denoted
$\pihat(X_i)$; for the simulation study of Section~\ref{sec:methodology}
we set $\pihat = \pi$.

\medskip
\noindent\textbf{Latent structural-zero indicator.}
In the zero-inflated regime, the response splits into a
structural-zero component and a count component. Let
$\Si \in \{0, 1\}$ be the latent indicator that unit $i$ is a
structural zero:
\begin{equation}
  \Si = 1 \;\Longrightarrow\; Y_i = 0, \qquad
  \Si = 0 \;\Longrightarrow\; Y_i \mid X_i, Z_i \sim \Poi(\lambda(X_i, Z_i))
                              \quad (\text{or } \NB(\lambda, \kappa)),
  \label{eq:Si}
\end{equation}
and write the conditional structural-zero probability as
\begin{equation}
  \pzi(X_i) \;\equiv\; \Pr(\Si = 1 \mid X_i, Z_i = z),
  \qquad z \in \{0, 1\}.
  \label{eq:pzi}
\end{equation}
The marginal count-component mean is
$\lambda(x, z) = \E[Y_i \mid X_i = x, Z_i = z, \Si = 0]$; we
parameterise $\log \lambda$ via a sum-of-trees prior in
Section~\ref{sec:countbcf}. The non-zero-inflated regime is recovered
by setting $\pzi \equiv 0$ (no structural-zero process); equivalently,
$\Si \equiv 0$ for all $i$.

\medskip
\noindent\textbf{Response-scale potential-outcome means.}
Combining the two channels yields the response-scale potential-outcome
means
\begin{equation}
  \muz(x) \;=\; \E[Y_i(z) \mid X_i = x]
          \;=\; \bigl(1 - \pzi(x)\bigr) \,\exp\bigl(\log \lambda_z(x)\bigr),
  \qquad z \in \{0, 1\}.
  \label{eq:potmeans}
\end{equation}
Our target estimands are the \emph{response-scale CATE} and ATE:
\begin{equation}
  \tauresp(x) \;=\; \muone(x) - \muzero(x), \qquad
  \mathrm{ATE} \;=\; \E_{X}\!\bigl[\tauresp(X)\bigr].
  \label{eq:cate-ate}
\end{equation}
When $\pzi \equiv 0$, \eqref{eq:potmeans} reduces to
$\muz(x) = \exp(\log\lambda_z(x))$ and \eqref{eq:cate-ate} becomes the
plain count CATE.

\medskip
\noindent\textbf{Two-channel decomposition.}
In the zero-inflated regime, the response-scale CATE in
\eqref{eq:cate-ate} mixes a change in the count intensity and a
change in the structural-zero probability. Methods that fit only one
channel necessarily miss part of the effect. In
Section~\ref{sec:summaries} we therefore report -- alongside
$\tauresp$ -- the two per-unit channel decompositions
\begin{equation}
  \Deltalam(x) \;=\; \exp(\log \lambda_1(x)) - \exp(\log \lambda_0(x)),
  \qquad
  \Deltapzi(x) \;=\; \pzi[1](x) - \pzi[0](x).
  \label{eq:channels}
\end{equation}

\medskip
\noindent\textbf{Scales.}
Treatment effects naturally live on up to three scales:
\begin{itemize}
  \item the \emph{count log-rate scale}, on which the moderating forest
        $\tauF(x)$ is parameterised; this is a relative-risk-like
        quantity, $\tauF(x) \approx \log\{\lambda_1(x) / \lambda_0(x)\}$
        in the non-zero-inflated case;
  \item the \emph{ZI log-odds scale}, on which the treatment shift in
        $\etazi$ is parameterised (zero-inflated regime only);
  \item the \emph{response (rate) scale}, $\tauresp(x)$ in
        \eqref{eq:cate-ate}, against which all reported metrics are
        evaluated.
\end{itemize}
All relevant quantities are reported. The simulation study
(Section~\ref{sec:methodology}) is scored against the rate-scale
truth.

% =======================================================================
\section{Background: log-linear BART}\label{sec:background}

This section gives the condensed restatement of the log-linear BART
construction of \citet{murray2021loglinear} -- including the
zero-inflation extension of \S4.2.2 -- that we need to define
CountBCF in Section~\ref{sec:countbcf}. The prognostic and moderating
forests of each function group are individually log-linear BART
regressions, so the leaf prior, the conditional sufficient
statistics, and the latent-variable augmentation recapped here
reappear unchanged inside the CountBCF Gibbs sampler. A reader
already fluent in \citet{murray2021loglinear} \S\S3--4.2.2 may skip to
Section~\ref{sec:countbcf}.

\subsection{The BART prior}\label{sec:bart-prior}

A single BART regression \citep{chipman2010bart} expresses an unknown
function $f : \mathbb{R}^p \to \mathbb{R}$ as a sum of $m$ shallow
regression trees,
\begin{equation}
  f(\x) \;=\; \sum_{h=1}^{m} g(\x;\; T_h, M_h),
  \label{eq:bart-sum}
\end{equation}
where $T_h$ is a binary tree partitioning $\mathbb{R}^p$ and
$M_h = (\mu_{h,1}, \dots, \mu_{h, b_h})'$ collects the leaf parameters
at the $b_h$ terminal nodes of $T_h$; the operator $g(\x;\; T_h, M_h)$
returns $\mu_{h,t}$ when $\x$ falls in terminal node $t$ of $T_h$. The
tree prior of \citet{chipman1998bayesian} assigns each internal node
at depth $d$ a split probability $\alpha\,(1+d)^{-\beta}$, with
$(\alpha, \beta) = (0.95, 2)$ biasing trees toward stubs. With a
Gaussian response, a Gaussian leaf prior is conjugate, and the
posterior is sampled by a backfitting Gibbs scheme that updates one
tree at a time using grow / prune / change / swap MH moves on $T_h$
and a conjugate draw of $M_h \mid T_h$.

\subsection{The log-linear extension}\label{sec:loglinear-extension}

For non-negative responses one writes the same sum-of-trees
representation on the log scale,
\begin{equation}
  \log f(\x) \;=\; \sum_{h=1}^{m} g(\x;\; T_h, M_h),
  \label{eq:loglinear-bart}
\end{equation}
and reparameterises each leaf parameter as
$\lambda_{h,t} = \exp(\mu_{h,t}) > 0$, so that
$f(\x) = \prod_h \lambda_{h, t(\x; T_h)}$. The Poisson and negative
binomial regression models take the form
\begin{equation}
  Y_i \mid \X_i \;\sim\; \Poi\bigl(\mu_{0i}\, f(\X_i)\bigr),
  \qquad
  Y_i \mid \X_i \;\sim\; \NB\bigl(\mathrm{size} = \kappa,\;
                                  \mathrm{mean} = \mu_{0i}\, f(\X_i)\bigr),
  \label{eq:loglinear-likelihoods}
\end{equation}
with $\mu_{0i}$ a unit-specific offset and $\kappa$ the NB dispersion,
satisfying $\Var(Y_i) = \mu_{0i} f(\X_i) + \mu_{0i}^2 f(\X_i)^2 /
\kappa$. The Poisson model is the $\kappa \to \infty$ limit. Both
likelihoods reappear inside the zero-inflated construction of
Section~\ref{sec:zi-param}: the not-structural-zero component of the
ZI mixture is exactly the model in \eqref{eq:loglinear-likelihoods},
while the structural-zero process is governed by an independent pair
of log-linear BART components introduced below.

After collecting all but the $h$th tree's fit into
$f^{-h}(\X_i) = \prod_{l \neq h} g(\X_i; T_l, \Lambda_l)$, the
contribution of a single observation to the joint posterior over
$(T_h, \Lambda_h)$ has the conditionally Gamma-shaped form
\begin{equation}
  L\bigl((T_h, \Lambda_h);\; T_{-h}, \Lambda_{-h}, y\bigr)
  \;\propto\;
  \prod_{t=1}^{b_h}
    \lambda_{h,t}^{r_{h,t}}\,
    \exp\!\bigl[-s_{h,t}\, \lambda_{h,t}\bigr],
  \label{eq:bart-poisson-leaflik}
\end{equation}
with conditional sufficient statistics
\begin{equation}
  r_{h,t} \;=\; \sum_{i:\, \X_i \in \mathcal{A}_{h,t}} u_i,
  \qquad
  s_{h,t} \;=\; \sum_{i:\, \X_i \in \mathcal{A}_{h,t}} v_i\, f^{-h}(\X_i),
  \label{eq:bart-suffstats}
\end{equation}
$\mathcal{A}_{h,t}$ the leaf-$t$ partition cell of $T_h$, and
$(u_i, v_i)$ likelihood-specific summaries that depend on the function
group (see Sections~\ref{sec:loglinear-augmentation} and
\ref{sec:zi-param}).

\subsection{The GIG-mixture leaf prior}\label{sec:gig-leaf-prior}

\citet{murray2021loglinear} constructs a leaf prior conjugate to
\eqref{eq:bart-poisson-leaflik} from two pieces. The generalised
inverse Gaussian (GIG) family has density
\begin{equation}
  p_{\GIG}(\lambda \mid \eta, \chi, \psi)
  \;\propto\;
  \lambda^{\eta - 1}\,
  \exp\!\Bigl[-\tfrac{1}{2}\bigl(\chi / \lambda + \psi \lambda\bigr)\Bigr],
  \qquad \lambda > 0,
  \label{eq:gig-density}
\end{equation}
covering the Gamma ($\chi = 0$) and inverse-Gamma ($\psi = 0$)
families as boundary cases. The CountBCF leaf prior is the equal
mixture
\begin{equation}
  p\bigl(\lambda_{h,t} \mid \leafc, \leafd\bigr)
  \;=\;
  \tfrac{1}{2}\, p_{\GIG}\!\bigl(\lambda_{h,t} \mid -\leafc,\, 2\leafd,\, 0\bigr)
  \;+\;
  \tfrac{1}{2}\, p_{\GIG}\!\bigl(\lambda_{h,t} \mid +\leafc,\, 0,\, 2\leafd\bigr),
  \label{eq:gig-leaf-prior}
\end{equation}
with stochastic representation $\lambda_{h,t} = W$ with probability
$1/2$ and $\lambda_{h,t} = 1/W$ otherwise, $W \sim \Ga(\leafc,
\leafd)$. This construction makes $\log \lambda_{h,t}$ exactly
symmetric about zero and keeps \eqref{eq:gig-leaf-prior} conjugate to
\eqref{eq:bart-poisson-leaflik}: the leaf full conditional is again a
GIG mixture, and the integrated likelihood for the tree MH step is a
ratio of GIG normalising constants \citep[\S3.1]{murray2021loglinear}.

The pair $(\leafc, \leafd)$ is reparameterised by a user-facing
concentration -- $\azero$ for the count component, $\zconc$ for the
ZI components (Section~\ref{sec:zi-param}) -- and the number of trees
$m$, by setting $\psi_1(\leafc) = \azero^2 / m$ and
$\leafd = \exp(\psi(\leafc))$, where $\psi$ and $\psi_1$ are the
digamma and trigamma functions. This calibration gives
$\E[\log \lambda_{h,t}] = 0$ and
$\Var(\log \lambda_{h,t}) = \azero^2 / m$, so the implied prior on
$\log f(\X) = \sum_h \log \lambda_{h, t(\X; T_h)}$ is approximately
$\N(0, \azero^2)$ at any $\X$. The concentration therefore plays the
same role for $\log f$ that the prior leaf scale plays for $f$ in
Gaussian BART. For large $m$ or small $\azero$ the calibration admits
the closed-form approximation $\leafc \approx m / \azero^2 + 1/2$,
$\leafd \approx m / \azero^2$ \citep[Proposition
3]{murray2021loglinear}.

\subsection{Latent-variable augmentation for the count
component}\label{sec:loglinear-augmentation}

To put the count likelihood into the form
\eqref{eq:bart-poisson-leaflik}, \citet{murray2021loglinear}
introduces one Gamma latent variable $\xilat$ per observation, with
full conditional
$\xilat \mid - \sim \Ga(\kappa + y_i,\; \kappa + \mu_{0i} f(\X_i))$
under the NB likelihood, and $\xilat \equiv 1$ under the Poisson
likelihood. Conditioning on $\xilat$, the per-observation
contribution of $f$ collapses to a Poisson-type kernel
$f(\X_i)^{y_i}\, \exp[-\xilat \mu_{0i} f(\X_i)]$, which integrated
against the GIG-mixture prior \eqref{eq:gig-leaf-prior} gives
\eqref{eq:bart-poisson-leaflik} with $u_i = y_i$ and
$v_i = \xilat \mu_{0i}$. Under the NB likelihood the dispersion
$\kappa$ is updated by a random-walk MH step on $\log \kappa$ with a
$\Be$-prime prior implied by
$\kappa / (1 + \kappa) \sim \Be(\kappa_a, \kappa_b)$. We defer the
specific update to Section~\ref{sec:posterior}.

\subsection{Zero-inflation and the redundant
$(\fzeroSym, \foneSym)$ parameterisation}\label{sec:zi-param}

For a zero-inflated count outcome the likelihood is a two-component
mixture: with conditional probability $\omega(\X_i)$ the response is
a structural zero, and otherwise it is drawn from the Poisson or NB
law \eqref{eq:loglinear-likelihoods}. \citet{murray2021loglinear}
\S4.2.2 writes the mixing probability through a pair of log-linear
BART forests $\fzeroSym, \foneSym$,
\begin{equation}
  \omega(\X_i)
  \;=\;
  \frac{\exp(\fzeroSym(\X_i))}{\exp(\fzeroSym(\X_i)) + \exp(\foneSym(\X_i))}
  \;=\;
  \sigma\!\bigl(\fzeroSym(\X_i) - \foneSym(\X_i)\bigr),
  \label{eq:omega-redundant}
\end{equation}
where each of $\fzeroSym, \foneSym$ has its own log-linear BART prior
\eqref{eq:loglinear-bart} with an independent GIG-mixture leaf prior
\eqref{eq:gig-leaf-prior} (with concentration $\zconc$). The
parameterisation is \emph{redundant}: adding a common constant to
both $\fzeroSym$ and $\foneSym$ leaves $\omega$ unchanged. The
redundancy is identified through the prior -- both forests are
centred at zero with symmetric GIG-mixture leaf priors -- and confers
two concrete benefits. Conditionally Gamma-shaped leaf updates with
closed-form GIG-mixture posteriors are obtained via a pair of latent
variables. Let $\zilat \in \{0, 1\}$ be the augmented mixture
indicator that resolves whether an observed $Y_i = 0$ came from the
structural process ($\zilat = 0$) or from the count process with a
zero draw ($\zilat = 1$). When $Y_i > 0$, $\zilat \equiv 1$; when
$Y_i = 0$, $\zilat$ has the Bernoulli full conditional
\begin{equation}
  \Pr(\zilat = 1 \mid -)
  \;=\;
  \frac{(1 - \omega(\X_i))\, p_{\mathrm{count}}(0 \mid \X_i)}
       {\omega(\X_i) + (1 - \omega(\X_i))\, p_{\mathrm{count}}(0 \mid \X_i)},
  \label{eq:zilat-fc}
\end{equation}
where $p_{\mathrm{count}}(0 \mid \X_i)$ is the count-component zero
probability ($\exp(-\lambda)$ for Poisson / ZIP,
$(\kappa / (\kappa + \lambda))^\kappa$ for NB / ZINB). Let further
$\philat > 0$ be an Exponential latent decoupling $\fzeroSym$ and
$\foneSym$, with full conditional
\begin{equation}
  \philat \;\mid\; - \;\sim\;
  \Exp\!\bigl(\exp(\fzeroSym(\X_i)) + \exp(\foneSym(\X_i))\bigr).
  \label{eq:phi-fc}
\end{equation}
Conditioning on $(\zilat, \philat)$, the joint conditional likelihood
of $(\fzeroSym, \foneSym)$ factorises into independent terms of the
form \eqref{eq:bart-poisson-leaflik}: the $\fzeroSym$ leaves have
$(u_i, v_i) = (1 - \zilat, \philat)$ and the $\foneSym$ leaves have
$(u_i, v_i) = (\zilat, \philat)$ \citep[\S A.5, Proposition
1]{murray2021loglinear}. The redundant sampler is also numerically
better behaved than a strictly identified parameterisation:
\citet[Supplement \S A]{murray2021loglinear} report substantially
better effective-sample-size performance across the full
probability range. In the body of this paper we write the
structural-zero log-odds as $\etazi = \fzeroSym - \foneSym$ and
report inference for the identified functionals only.

\subsection{Implementation backbone}\label{sec:backend-note}

We inherit the C++ implementation of the log-linear BART sampler --
the GIG-mixture leaf draw, the grow / prune / change / swap tree MH
moves, the latent-variable updates $(\zilat, \xilat, \philat)$, the
soft-mixture-weight recomputation, and the dispersion MH step -- from
the \texttt{countbart} R package of \citet{wikle2023causal}. CountBCF
adds the BCF $(\muF, \tauF)$ decomposition on top of this backend by
replacing the single forest of log-linear BART (per function group)
with a forest pair and introducing a per-observation weight that
routes treated and control units to the correct forest; we describe
that change next.

% =======================================================================
\section{CountBCF: BCF decomposition for count and ZI count
outcomes}\label{sec:countbcf}

This section defines CountBCF. Section~\ref{sec:likelihood} states
the likelihood and parameterisation in both the count and the
zero-inflated regimes; Section~\ref{sec:priors} collects the priors
and the differential shrinkage applied symmetrically across function
groups; Section~\ref{sec:posterior} writes out the MCMC sweep in
Algorithm~\ref{alg:countbcf} and the associated Gibbs and MH updates;
Section~\ref{sec:summaries} explains how posterior draws are
post-processed into the estimands of Section~\ref{sec:notation},
including the per-unit channel decomposition in the ZI case.

\subsection{Likelihood and parameterisation}\label{sec:likelihood}

CountBCF models the conditional mean of a count response (and, in the
ZI regime, the structural-zero probability) on the log scale using a
treatment-aware BCF decomposition. We treat the count and
zero-inflated cases as two sub-models of a common framework, selected
by the user via a \texttt{count\_model} switch that takes values in
$\{\texttt{poisson}, \texttt{nb}, \texttt{zipoisson}, \texttt{zinb}\}$.

\paragraph{Count regime: three forests.}
For unit $i$ with covariates $X_i$ and treatment indicator
$Z_i \in \{0, 1\}$, the Poisson / NB sub-model writes
\begin{equation}
  \boxed{\;
    \log \lambda(X_i, Z_i)
    \;=\;
    \muF\!\bigl(X_i,\; \pihat(X_i)\bigr)
    \;+\;
    Z_i \cdot \tauF(X_i),
  \;}
  \label{eq:countbcf-loglambda}
\end{equation}
with response
\begin{equation}
  Y_i \mid X_i, Z_i
  \;\sim\;
  \Poi\bigl(\lambda(X_i, Z_i)\bigr)
  \;\;\text{or}\;\;
  \NB\bigl(\mathrm{size} = \kappa,\; \mathrm{mean} = \lambda(X_i, Z_i)\bigr),
  \label{eq:countbcf-likelihoods}
\end{equation}
where the size parameter $\kappa > 0$ governs overdispersion. The
Poisson model is the $\kappa \to \infty$ limit. Decomposition
\eqref{eq:countbcf-loglambda} is the count analogue of the additive
BCF decomposition $\E[Y \mid X, Z] = \muF(X, \pihat) + Z \tauF(X)$ of
\citet{hahn2020bayesian}. Both pieces are log-linear BART forests in
the sense of Section~\ref{sec:loglinear-extension}: $\muF$ is a sum
of $m_\mu$ trees with leaf parameters drawn from the GIG mixture
\eqref{eq:gig-leaf-prior}, and similarly $\tauF$ is a sum of $m_\tau$
trees. The prognostic forest models the control-arm log-mean across
the whole sample; the moderating forest enters only for treated units
and captures treatment-effect heterogeneity. As in BCF, $\pihat$ is
appended to the covariates feeding $\muF$ but \emph{not} those
feeding $\tauF$.

\paragraph{Zero-inflated regime: six forests, three groups.}
The ZIP / ZINB sub-model attaches a BCF $(\mu, \tau)$ decomposition
to each of the three log-linear BART forests $f, \fzeroSym, \foneSym$
of Section~\ref{sec:zi-param}. We index the three function groups by
$g \in \{0, 1, 2\}$ with $g = 0$ the count log-rate, $g = 1$ the
structural-zero log-odds component $\fzeroSym$, and $g = 2$ the
not-structural-zero log-odds component $\foneSym$. Then
\begin{equation}
  \boxed{\;
    \begin{aligned}
      \log \lambda(X_i, Z_i)
        &\;=\; \muF\!\bigl(X_i,\, \pihat(X_i)\bigr)
               \;+\; Z_i \cdot \tauF(X_i), \\[2pt]
      \etazi(X_i, Z_i)
        &\;=\; \bigl[\muFzero\!\bigl(X_i,\, \pihat(X_i)\bigr)
                     + Z_i \cdot \tauFzero(X_i)\bigr] \\
        &\quad - \bigl[\muFone\!\bigl(X_i,\, \pihat(X_i)\bigr)
                       + Z_i \cdot \tauFone(X_i)\bigr], \\[2pt]
      \pzi(X_i)
        &\;=\; \sigma\!\bigl(\etazi(X_i,\, Z_i = z)\bigr), \qquad z \in \{0, 1\},
    \end{aligned}
  \;}
  \label{eq:zicountbcf-loglambda}
\end{equation}
together with the ZI count likelihood
$Y_i \mid \Si = 1 = 0$, and
$Y_i \mid \Si = 0, X_i, Z_i \sim \Poi(\lambda(X_i, Z_i))$ or
$\NB(\kappa, \lambda(X_i, Z_i))$, where
$\Pr(\Si = 1 \mid X_i, Z_i) = \pzi(X_i)$. The non-ZI sub-model
collapses \eqref{eq:zicountbcf-loglambda} by dropping the $\etazi$
line and the two ZI groups, recovering \eqref{eq:countbcf-loglambda}.
Equation \eqref{eq:zicountbcf-loglambda} contains six unknown
functions, two per function group, summarised in
Table~\ref{tab:zicountbcf-forests}. Each is itself a log-linear BART
regression: a sum of $m_{s, g}$ shallow trees with leaf parameters
drawn from the GIG mixture, with $s \in \{\mu, \tau\}$ indexing the
BCF role and $g \in \{0, 1, 2\}$ indexing the function group. For
each group, the prognostic forest $\mu^{(g)}$ enters with weight
$\omegai = 1$; the moderating forest $\tau^{(g)}$ enters with weight
$\omegai = Z_i$. The propensity score estimate $\pihat$ is appended
to the covariates of all three prognostic forests but not to the
moderating ones.

\begin{table}[t]
\centering
\caption{The six forests of CountBCF in the zero-inflated regime. The
function group $g$ labels the log-linear BART target inherited from
\citet{murray2021loglinear}; the BCF role $s$ labels the prognostic
($\mu$) vs.\ moderating ($\tau$) piece introduced by the
$(\mu, \tau)$ decomposition. The non-ZI sub-model retains only the
$g = 0$ row.}
\label{tab:zicountbcf-forests}
\begin{tabular}{cllcc}
\toprule
$g$ & Function group & Forest pair $(\mu^{(g)}, \tau^{(g)})$
                                                      & Trees $(m_\mu, m_\tau)$ & Weight $\omegai$ on $\tau^{(g)}$ \\
\midrule
$0$ & Count log-rate $\log\lambda$           & $(\muF,\; \tauF)$
                                                      & $(250,\; 50)$ & $Z_i$ \\
$1$ & ZI log-odds $\fzeroSym$                & $(\muFzero,\; \tauFzero)$
                                                      & $(100,\; 50)$ & $Z_i$ \\
$2$ & not-ZI log-odds $\foneSym$             & $(\muFone,\; \tauFone)$
                                                      & $(100,\; 50)$ & $Z_i$ \\
\bottomrule
\end{tabular}
\end{table}

\medskip
\noindent\textbf{Scales recovered.}
On the \emph{count log-rate scale}, $\tauF(x)$ is itself a CATE: the
conditional log relative rate (of the not-structural-zero
sub-population in the ZI case),
$\tauF(x) = \log\{\lambda_1(x) / \lambda_0(x)\}$. On the \emph{ZI
log-odds scale}, the treatment shift in the structural-zero log-odds
is
\begin{equation}
  \Detazi(x)
  \;\equiv\;
  \etazi(x, 1) - \etazi(x, 0)
  \;=\;
  \tauFzero(x) - \tauFone(x),
  \label{eq:Deltaeta-decomp}
\end{equation}
so $\tauFzero, \tauFone$ are jointly identified up to a common shift
-- only the difference is statistically relevant for the ZI channel.
On the \emph{response (rate) scale}, the potential-outcome means
\eqref{eq:potmeans} are
\begin{equation}
  \muz(x)
  \;=\;
  \bigl(1 - \sigma(\etazi(x, z))\bigr)\,
  \exp\!\bigl(\muF(x, \pihat(x)) + z\, \tauF(x)\bigr),
  \qquad z \in \{0, 1\},
  \label{eq:muz-decomp}
\end{equation}
and $\tauresp(x) = \muone(x) - \muzero(x)$. In the non-ZI sub-model
$\pzi \equiv 0$ and \eqref{eq:muz-decomp} reduces to
$\muz(x) = \exp(\muF + z\tauF)$, giving the closed-form count CATE
\begin{equation}
  \tauresp(x)
  \;=\;
  \exp\!\bigl(\muF(x, \pihat(x))\bigr)\,
  \bigl(\exp(\tauF(x)) - 1\bigr),
  \label{eq:tauresp-decomp}
\end{equation}
which makes explicit that the absolute treatment effect on the rate
scale depends on the baseline rate in addition to the log-rate
effect. In the ZI case, $\tauresp$ also mixes the count-rate channel
$\Deltalam$ and the ZI-probability channel $\Deltapzi$ from
\eqref{eq:channels}.

\medskip
\noindent\textbf{Why this parameterisation.}
Three properties of the construction are worth flagging. First, the
model respects the non-negativity of $Y$ automatically: $\lambda > 0$
and (where applicable) $\pzi \in (0, 1)$ regardless of the forest
values, so no truncation or clipping is required. Second, all forests
are parameterised \emph{a priori} independently, which is what
permits the BCF differential-shrinkage idea
\citep[\S2.3]{hahn2020bayesian} to be applied per group: each
$\tau^{(g)}$ can be shrunk harder than its $\mu^{(g)}$ counterpart
without disturbing other groups' priors. Third, the construction
reduces to log-linear BART \citep{murray2021loglinear} when
$\tau^{(g)} \equiv 0$ for all $g$, and to standard BCF in the limit
where $\tauF$ is shrunk hard to zero. In this sense CountBCF is the
smallest extension of log-linear BART that delivers a separate
moderating forest on each channel, and the smallest extension of BCF
that handles count and zero-inflated count outcomes natively. The
redundant $(\fzeroSym, \foneSym)$ parameterisation in the ZI case is
retained for the reasons given in Section~\ref{sec:zi-param}:
conjugate leaf updates, well-defined posterior on the identified
functionals, and substantially better MCMC mixing
\citep[Supplement~\S A]{murray2021loglinear}. The BCF $(\mu, \tau)$
decomposition does not break the redundancy: shifting $\muFzero$ and
$\muFone$ by a common constant leaves $\etazi$ invariant, and shifting
$\tauFzero$ and $\tauFone$ by a common constant leaves $\Detazi$
invariant.

\subsection{Priors}\label{sec:priors}

The CountBCF prior factorises across forests, the dispersion
parameter, and the GIG-mixture latent structure of each leaf. The
construction is symmetric across function groups, so we state the
prior for a generic forest pair $(\mu^{(g)}, \tau^{(g)})$ and
specialise the concentration hyperparameters to the count and ZI
groups.

\paragraph{Sum-of-trees priors on each forest.}
Each forest is a log-linear BART regression in the sense of
Section~\ref{sec:loglinear-extension}: a sum of $m_{s,g}$ shallow
trees with tree structures $\{T_{s,g,h}\}$ drawn \emph{a priori}
independently from the CGM tree-structure prior and leaf parameters
$\{\Lambda_{s,g,h}\} = \{\lambda_{s, g, h, t}\}$ drawn \emph{a priori}
independently from the GIG mixture \eqref{eq:gig-leaf-prior}.
Collecting these, for $s \in \{\mu, \tau\}$ and (in the ZI case)
$g \in \{0, 1, 2\}$,
\begin{equation}
  \boxed{\;
    \begin{aligned}
      T_{s, g, h} \;\big|\; \alpha_s, \beta_s
        &\;\stackrel{\mathrm{iid}}{\sim}\;
          \mathrm{CGM}(\alpha_s, \beta_s),
          & h &= 1, \dots, m_{s, g}, \\[2pt]
      \lambda_{s, g, h, t} \;\big|\; \leafc_{s,g}, \leafd_{s,g}
        &\;\stackrel{\mathrm{iid}}{\sim}\;
          \tfrac{1}{2}\,\GIG(-\leafc_{s,g}, 2\leafd_{s,g}, 0)
          + \tfrac{1}{2}\,\GIG(\leafc_{s,g}, 0, 2\leafd_{s,g}),
          & t &= 1, \dots, b_{s, g, h}, \\[2pt]
      \mu^{(g)}(x, \pihat(x))
        &\;=\; \sum_{h=1}^{m_{\mu, g}} g\!\bigl((x, \pihat(x));\, T_{\mu, g, h}, \Lambda_{\mu, g, h}\bigr),
        & \tau^{(g)}(x)
        &\;=\; \sum_{h=1}^{m_{\tau, g}} g\bigl(x;\, T_{\tau, g, h}, \Lambda_{\tau, g, h}\bigr),
    \end{aligned}
  \;}
  \label{eq:zicountbcf-prior}
\end{equation}
where the CGM density assigns an internal-node split probability
$\alpha_s\,(1+d)^{-\beta_s}$ at depth $d$ and a uniform draw over
splitting variable and threshold given the split. The
depth-penalty hyperparameters $(\alpha_s, \beta_s)$ are shared across
function groups (so all prognostic forests have the same tree-shape
prior, and similarly for the moderating forests); the leaf
concentration is allowed to differ across $g$ to reflect the
different natural scales of the count log-rate versus the
zero-inflation log-odds. In the count-only regime, indices $g \in
\{1, 2\}$ are dropped throughout.

\paragraph{Differential shrinkage applied symmetrically.}
CountBCF inherits the BCF \citep{hahn2020bayesian} idea that the
treatment-effect surface $\tau^{(g)}$ should be shrunk more
aggressively than the prognostic surface $\mu^{(g)}$ on \emph{each}
channel: the signal-to-noise in $\tau^{(g)}$ is typically lower than
in $\mu^{(g)}$; spurious heterogeneity in any of $\tauF, \tauFzero,
\tauFone$ manifests as fake interactions with $Z$ on the respective
channel; and a stub-biased prior on every $\tau^{(g)}$ delivers a
near-constant treatment-effect estimator as a default. The default
hyperparameters are summarised in
Table~\ref{tab:default-hyperparams}. The tree-prior
$(\alpha_\tau, \beta_\tau) = (0.25, 3)$ sharply downweights non-stub
trees in every moderating forest: under it the prior probability of a
non-trivial split at the root is $0.25$ and the probability of a
depth-$2$ internal node is approximately $0.0039$, so on average each
moderating forest is close to a piecewise-constant surface with very
few splits. The leaf concentration $\azerotau = \azero / 2$ (count
group) and $\zconctau = \zconc / 2$ (ZI groups) further shrinks the
amplitude of those few splits: under the GIG-mixture calibration of
Section~\ref{sec:gig-leaf-prior},
$\Var(\log \lambda_{\tau, g, h, t}) \approx a_{0, \tau, g}^2 /
m_{\tau, g}$, so the implied marginal prior on $\tau^{(g)}(x)$ is
$\N(0, a_{0, \tau, g}^2)$, a factor of two tighter than the prior on
$\mu^{(g)}$.

The user-facing concentration $\azero$ defaults to
$\tfrac{1}{2}(\log y_{(0.95)} - \overline{\log \mu_{0,i}})$, where
$y_{(0.95)}$ is the empirical 95\% quantile of $y$, centring the
count-channel prior on a data-matched scale. The ZI concentration
$\zconc$ defaults to $3.5 / \sqrt{2} \approx 2.47$, matching the
calibration of \citet[\S 4.2.2]{murray2021loglinear} for the
redundant $(\fzeroSym, \foneSym)$ parameterisation. Both can be
overridden by the user.

\begin{table}[t]
\centering
\caption{CountBCF default hyperparameters. The ZI rows ($g \in
\{1, 2\}$) apply only when \texttt{count\_model} $\in
\{\texttt{zipoisson}, \texttt{zinb}\}$; the dispersion parameters
apply only when \texttt{count\_model} $\in \{\texttt{nb},
\texttt{zinb}\}$.}
\label{tab:default-hyperparams}
\begin{tabular}{lll}
\toprule
Hyperparameter & Default & Role \\
\midrule
$m_{\mu, 0}$, $m_{\tau, 0}$ & $250$,\, $50$ & Trees per forest, count group $g = 0$ \\
$m_{\mu, 1}$, $m_{\tau, 1}$ & $100$,\, $50$ & Trees per forest, ZI group $g = 1$ ($\fzeroSym$) \\
$m_{\mu, 2}$, $m_{\tau, 2}$ & $100$,\, $50$ & Trees per forest, ZI group $g = 2$ ($\foneSym$) \\
$(\alpha_\mu, \beta_\mu)$ & $(0.95,\, 2)$ & CGM depth penalty, all $\mu^{(g)}$ \\
$(\alpha_\tau, \beta_\tau)$ & $(0.25,\, 3)$ & CGM depth penalty, all $\tau^{(g)}$ \\
$\azero$ & $\tfrac{1}{2}(\log y_{(0.95)} - \overline{\log \mu_{0,i}})$ & Leaf concentration, $\muF$ (count) \\
$\azerotau$ & $\azero / 2$ & Leaf concentration, $\tauF$ (count) \\
$\zconc$ & $3.5 / \sqrt{2}$ & Leaf concentration, $\muFzero, \muFone$ (ZI) \\
$\zconctau$ & $\zconc / 2$ & Leaf concentration, $\tauFzero, \tauFone$ (ZI) \\
$(\kappa_a, \kappa_b)$ & $(5,\, 3)$ & Beta-prime prior on $\kappa$ \\
$\sigma_\kappa$ & $0.21$ & RW-MH proposal SD on $\log \kappa$ \\
\texttt{include\_pihat} & \texttt{"control"} & Propensity appended to all $\mu^{(g)}$ \\
\bottomrule
\end{tabular}
\end{table}

\paragraph{Propensity-as-covariate fix on every channel.}
A central insight of \citet{hahn2020bayesian} is that, on the
additive response scale, naively regressing the outcome on $(X, Z)$
with a BART-like prior leaves the estimator vulnerable to
regularisation-induced confounding (RIC): when the propensity score
$\pi(x)$ is a function of $x$, the prior shrinkage of the prognostic
surface toward zero biases $\tauF$ in a direction that depends on
$\pi$. The fix is to append $\pihat(x)$ to the prognostic forest's
covariates so that $\muF$ can ``absorb'' the part of the outcome
explained by $\pi$ without recruiting $\tauF$. On the log-rate scale
of \eqref{eq:countbcf-loglambda} the RIC argument is at least as
severe: a perturbation $\delta$ in $\muF$ translates to a
multiplicative perturbation $e^\delta - 1 \approx \delta$ in the
baseline rate, which then multiplies into $\tauresp$ via
\eqref{eq:tauresp-decomp}; a bias in $\muF$ correlated with $\pi$
propagates linearly into $\tauresp$. In the ZI case the same RIC
argument applies \emph{independently} on each channel: perturbations
in $\muFzero$ or $\muFone$ propagate into the ZI-channel moderating
forests through the redundant parameterisation. We therefore append
$\pihat$ to the covariate set of \emph{every} prognostic forest by
default and do \emph{not} append it to any moderating forest. The R
wrapper exposes this as \verb|include_pihat = "control"|, with
\verb|"moderate"|, \verb|"both"|, and \verb|"none"| available for
sensitivity checks.

\paragraph{Negative-binomial dispersion $\kappa$.}
Under the NB / ZINB likelihood, the size parameter $\kappa$ controls
overdispersion: $\Var(Y_i \mid X_i, Z_i, \Si = 0) = \lambda_i +
\lambda_i^2 / \kappa$. Following \citet{murray2021loglinear} we place
a Beta-prime prior on $\kappa$ implied by
\begin{equation}
  \frac{\kappa}{1 + \kappa} \;\sim\; \Be(\kappa_a, \kappa_b),
  \qquad
  (\kappa_a, \kappa_b) = (5,\, 3),
  \label{eq:kappa-prior}
\end{equation}
with mean $2.5$ and a right tail heavy enough not to rule out
near-Poisson behaviour while still penalising arbitrarily large
overdispersion. The same prior governs $\kappa$ under both NB and
ZINB; structural zeros enter the $\kappa$-conditional likelihood only
through $\zilat$ (Section~\ref{sec:posterior}). The dispersion is
updated by a random-walk MH step on $\log \kappa$ with Gaussian
proposal of standard deviation $\sigma_\kappa = 0.21$.

\subsection{Posterior computation}\label{sec:posterior}

Posterior inference for CountBCF uses a Gibbs sampler that
interleaves backfitting updates of the forests with latent-variable
augmentation steps and (under NB / ZINB) a MH step on the dispersion.
The structure mirrors the log-linear BART sampler of
\citet[\S 3.3, \S S.3]{murray2021loglinear}: the only forest-level
change is the introduction of a per-observation weight
$\omegai \in \{1, Z_i\}$ that routes treated and control units to the
moderating and prognostic forests, replicated across all function
groups.

\paragraph{Latent-variable augmentation.}
The augmentations are inherited from
Sections~\ref{sec:loglinear-augmentation} and \ref{sec:zi-param}.
Each observation carries up to three latents:
\begin{equation}
  \begin{aligned}
    \zilat \;\mid\; -
      &\;\sim\;
      \Ber\!\Bigl(
        \frac{(1 - \omega(\X_i))\, p_{\mathrm{count}}(0 \mid \X_i)}
             {\omega(\X_i) + (1 - \omega(\X_i))\,
              p_{\mathrm{count}}(0 \mid \X_i)}\Bigr)
        \quad\text{when } Y_i = 0 \text{ (ZIP/ZINB)};
        \;\; \zilat \equiv 1 \text{ when } Y_i > 0, \\[2pt]
    \xilat \;\mid\; -
      &\;\sim\;
      \begin{cases}
        \delta_{1} & \text{(Poisson / ZIP)}, \\
        \Ga\!\bigl(\kappa + Y_i,\;\kappa + \lambda(\X_i, Z_i)\bigr)
        & \text{(NB / ZINB)},
      \end{cases} \\[2pt]
    \philat \;\mid\; -
      &\;\sim\;
      \Exp\!\bigl(\exp(\fzeroSym(\X_i)) + \exp(\foneSym(\X_i))\bigr)
      \quad\text{(ZIP/ZINB only)}.
  \end{aligned}
  \label{eq:three-latents}
\end{equation}
Under the non-ZI sub-models the $\zilat$ and $\philat$ updates are
skipped; under the Poisson / ZIP sub-models the $\xilat$ update is
skipped.

\paragraph{Tree updates: a single primitive applied to all forests.}
Conditional on the latents, the joint conditional likelihood of the
forests factorises into independent terms, each of which has the
conditionally Gamma form \eqref{eq:bart-poisson-leaflik}. For forest
$s$ in function group $g$ at tree $h$, the conditional augmented
likelihood for the $b_{s,g,h}$ leaves of $T_{s,g,h}$ is
\begin{equation}
  \boxed{\;
    L\bigl(\Lambda_{s,g,h} \,\big|\, T_{s,g,h}, T_{-(s,g,h)},
           \Lambda_{-(s,g,h)},
           \zilat, \xilat, \philat, Y, \kappa\bigr)
    \;\propto\;
    \prod_{t=1}^{b_{s,g,h}}
      \lambda_{s,g,h,t}^{\,r_{s,g,h,t}}\;
      \exp\!\bigl[-s_{s,g,h,t}\,\lambda_{s,g,h,t}\bigr],
  \;}
  \label{eq:zi-augmented-leaflik}
\end{equation}
where the leaf sufficient statistics
$r_{s,g,h,t} = \sum_{i \in \mathcal{A}_{s,g,h,t}} u_i$ and
$s_{s,g,h,t} = \sum_{i \in \mathcal{A}_{s,g,h,t}}
v_i\, f^{-(s,g,h)}(\X_i)$ aggregate over the units in the partition
cell, using the per-group $(u_i, v_i)$ enumerated in
Table~\ref{tab:suffstats}. The BCF weight $\omegai \in \{1, Z_i\}$
enters multiplicatively in both $u_i$ and $v_i$: for prognostic
forests every unit contributes; for moderating forests only treated
units contribute. This is the BCF moderating-forest ingredient
transposed to the count (or ZI count) likelihood, applied
symmetrically to all function groups present in the model.

\begin{table}[t]
\centering
\caption{Per-group sufficient statistics. The values $u_i, v_i$ are
inserted into \eqref{eq:zi-augmented-leaflik} through the leaf
accumulators $r_{s,g,h,t}, s_{s,g,h,t}$. The BCF weight is
$\omegai = 1$ for prognostic forests ($s = \mu$) and $\omegai = Z_i$
for moderating forests ($s = \tau$); it premultiplies $(u_i, v_i)$
before the leaf sum is formed. The non-ZI sub-model retains only the
$g = 0$ row.}
\label{tab:suffstats}
\begin{tabular}{clll}
\toprule
$g$ & Function target          & $u_i \cdot \omegai^{-1}$ & $v_i \cdot \omegai^{-1}$ \\
\midrule
$0$ & count rate $\lambda$     & $\zilat\, y_i$
                                & $\zilat\, \exp(\log \xilat + \log\lambda(\X_i, Z_i)
                                            - \omegai \log \lambda_{s,g,h,t})$ \\
$1$ & ZI log-odds $\fzeroSym$  & $1 - \zilat$
                                & $\exp(\log \philat + \log\fzeroSym(\X_i)
                                            - \omegai \log \lambda_{s,g,h,t})$ \\
$2$ & ZI log-odds $\foneSym$   & $\zilat$
                                & $\exp(\log \philat + \log\foneSym(\X_i)
                                            - \omegai \log \lambda_{s,g,h,t})$ \\
\bottomrule
\end{tabular}
\end{table}

The GIG-mixture leaf prior \eqref{eq:gig-leaf-prior} is conjugate to
\eqref{eq:zi-augmented-leaflik}, so the leaf draw is a closed-form
mixture of GIG variates (implemented in \texttt{drmu\_loglinear}), and
$T_{s,g,h}$ is updated by a MH step with grow / prune / change / swap
proposals (implemented in \texttt{bd\_loglinear}). The MH acceptance
ratio uses the integrated marginal likelihood obtained by integrating
$\Lambda_{s,g,h}$ against the GIG-mixture prior; because both prior
and likelihood are products over leaves, that integral factorises and
reduces to a ratio of GIG normalising constants
\citep[Appendix~A]{murray2021loglinear}.

\paragraph{Per-group log-fit accumulator.}
Implemented as backfitting, the tree update maintains running
accumulators on the $n$-vector index $k$: $\mathtt{allfit}[k]$ tracks
the total log-mean / log-odds; $\mathtt{allfits}[s][k]$ tracks the
contribution of forest $s$; and the per-group accumulator
$\mathtt{log\_f\_per\_group}[g][k] = \mu^{(g)}(\X_k, \pihat(\X_k)) +
Z_k\, \tau^{(g)}(\X_k)$ stores the in-sample log-fit of both forests
in group $g$. Before sampling tree $(s, g, h)$ its current
contribution is subtracted from all relevant accumulators; sufficient
statistics are computed from the partial fit; $T_{s,g,h}$ and
$\Lambda_{s,g,h}$ are updated; and the new contribution is added back.
The per-group accumulator makes the six-forest update efficient: the
ZI sufficient statistics read $\log \fzeroSym$ and $\log \foneSym$
from $\mathtt{log\_f\_per\_group}[1]$ and
$\mathtt{log\_f\_per\_group}[2]$ in $O(1)$ per observation. After
each sweep the soft mixture weight that enters the $\zilat$ update is
recomputed as $\log\softw(\X_i) = \log\foneSym(\X_i) -
\logsumexp(\log\fzeroSym(\X_i), \log\foneSym(\X_i))$, which avoids the
underflow of naive computation when one of the log-fits is large and
negative.

\paragraph{Dispersion update (NB / ZINB only).}
Propose $\kappa^\star = \exp(\log \kappa + \sigma_\kappa\,
\varepsilon)$ with $\varepsilon \sim \N(0, 1)$. Accept with
probability
\begin{equation}
  \alpha_\kappa
  \;=\;
  \min\!\Bigl\{1,\;
    \frac{p_{\NB}(Y \mid \zilat, \lambda, \kappa^\star)}{p_{\NB}(Y \mid \zilat, \lambda, \kappa)}
    \cdot
    \frac{p(\kappa^\star)}{p(\kappa)}
    \cdot
    \frac{\kappa^\star}{\kappa}
  \Bigr\},
  \label{eq:kappa-MH}
\end{equation}
where the likelihood factor is the augmented NB likelihood conditional
on the mixture indicator (so structural zeros do not enter the
dispersion update under ZINB), $p(\kappa)$ is the Beta-prime density
implied by \eqref{eq:kappa-prior}, and the Jacobian factor $\kappa^\star
/ \kappa$ comes from the log-scale proposal. Acceptance rates after
burn-in are returned as \texttt{kappa\_acpt} and serve as a coarse
convergence diagnostic; we target the $0.20$--$0.45$ band typical for
random-walk MH.

\begin{algorithm}[t]
\caption{CountBCF: one MCMC sweep. Indices $s \in \{\mu, \tau\}$
range over forests, $g \in \{0\}$ in the count sub-model and
$g \in \{0, 1, 2\}$ in the ZI sub-model, and $h = 1, \dots, m_{s, g}$
over the trees of forest $(s, g)$. The BCF weight is $\omegai = 1$
for $s = \mu$ and $\omegai = Z_i$ for $s = \tau$. Latent updates that
are not applicable to the active sub-model are skipped.}
\label{alg:countbcf}
\begin{algorithmic}[1]
  \Require Current state
  $\bigl(\{T_{s,g,h}, \Lambda_{s,g,h}\},\, \kappa,\,
         \{\zilat, \xilat, \philat\}\bigr)$.
  \Statex
  \State \textbf{Dispersion update (NB / ZINB only).}
         Propose $\kappa^\star = \exp(\log \kappa + \sigma_\kappa\,
         \varepsilon)$, $\varepsilon \sim \N(0, 1)$. Accept with
         probability $\alpha_\kappa$ from \eqref{eq:kappa-MH}; otherwise
         retain $\kappa$.
  \Statex
  \State \textbf{Augmentation update for $\zilat$ (ZIP / ZINB only).}
         For each $i$ with $Y_i = 0$, draw $\zilat$ from the Bernoulli
         full conditional in \eqref{eq:three-latents}; set $\zilat = 1$
         when $Y_i > 0$. [\texttt{drz\_loglinear}]
  \Statex
  \State \textbf{Augmentation update for $\xilat$ (NB / ZINB only).}
         For each $i$, draw $\xilat$ from the Gamma full conditional;
         set $\xilat \equiv 1$ otherwise.
         [\texttt{drxi\_loglinear}]
  \Statex
  \State \textbf{Augmentation update for $\philat$ (ZIP / ZINB only).}
         For each $i$, draw
         $\philat \sim \Exp(\exp(\fzeroSym(\X_i)) +
         \exp(\foneSym(\X_i)))$.
         [\texttt{drphi\_loglinear}]
  \Statex
  \State \textbf{Backfitting tree updates.}
         For each function group $g$ active in the sub-model, each
         forest $s \in \{\mu, \tau\}$, and each tree $h = 1, \dots,
         m_{s, g}$:
         \begin{enumerate}[label*=5.\arabic*, leftmargin=*, topsep=2pt, itemsep=2pt]
            \item Subtract the current contribution of
                  $(T_{s,g,h}, \Lambda_{s,g,h})$ from
                  $\mathtt{allfit}$, $\mathtt{allfits}[s]$, and
                  $\mathtt{log\_f\_per\_group}[g]$.
            \item Compute the leaf sufficient statistics
                  $(r_{s, g, h, t}, s_{s, g, h, t})$ from
                  \eqref{eq:zi-augmented-leaflik} using the per-group
                  $(u_i, v_i)$ of Table~\ref{tab:suffstats},
                  premultiplied by the weight $\omegai$.
            \item Propose a tree move (grow / prune / change / swap)
                  for $T_{s, g, h}$ and accept with probability from
                  the integrated marginal likelihood (GIG-mixture
                  conjugacy). [\texttt{bd\_loglinear}]
            \item Draw the leaf values
                  $\Lambda_{s, g, h} \mid T_{s, g, h}, \zilat, \xilat,
                  \philat, Y, \kappa$ from the GIG-mixture posterior
                  implied by \eqref{eq:zi-augmented-leaflik}.
                  [\texttt{drmu\_loglinear}]
            \item Add the new contribution back into the affected
                  accumulators.
         \end{enumerate}
  \Statex
  \State \textbf{Soft mixture weight (ZIP / ZINB only).}
         For each $i$, recompute
         $\log \softw(\X_i) = \log\foneSym(\X_i) -
         \logsumexp(\log\fzeroSym(\X_i), \log\foneSym(\X_i))$.
  \Statex
  \State \textbf{(Optional) Random-effects update.} If enabled, draw
         the random-intercept block from its Gibbs full conditional;
         otherwise skip.
  \Statex
  \State \Return Updated state
         $\bigl(\{T_{s,g,h}, \Lambda_{s,g,h}\},\, \kappa,\,
                \{\zilat, \xilat, \philat\}\bigr)$.
\end{algorithmic}
\end{algorithm}

\paragraph{Implementation lineage.}
The C++ primitives \texttt{bd\_loglinear}, \texttt{drmu\_loglinear},
\texttt{fit\_loglinear}, \texttt{drz\_loglinear},
\texttt{drxi\_loglinear}, \texttt{drphi\_loglinear}, and the
dispersion MH step are inherited verbatim from \texttt{countbart}
\citep{wikle2023causal}; the CountBCF contribution is the dispatch
logic in \texttt{src/countbcf.cpp} that (i) instantiates two forests
per function group rather than one, (ii) sets the per-forest weight
vector $\omegai$ to either $\mathbf{1}_n$ or $Z_{\mathrm{trt}}$, (iii)
applies $\omegai$ to the sufficient-statistic construction in
step~5.2, and (iv) maintains the per-group log-fit accumulator so
that the four ZI forests read $\fzeroSym$ and $\foneSym$ in $O(1)$
per observation. The change is mechanical at the kernel level but
yields a sampler that, to our knowledge, is functionally distinct
from any log-linear BART variant in the published literature: it
jointly estimates up to six log-rate / log-odds surfaces under
differential shrinkage, with each moderating forest seeing only the
treated subsample's residual variation.

\subsection{Posterior summaries}\label{sec:summaries}

Algorithm~\ref{alg:countbcf} produces a sequence of $B$ post-burn-in
posterior draws indexed by $b = 1, \dots, B$. The sampler stores, for
each draw, the per-unit values of the active forests on the original
(unsorted) input order:
\begin{equation}
  \begin{aligned}
    \muF^{(b)}(X_i) &= \texttt{mu\_f\_post}[b, i],
    & \tauF^{(b)}(X_i) &= \texttt{tau\_f\_post}[b, i], \\
    \muFzero^{(b)}(X_i) &= \texttt{mu\_f0\_post}[b, i],
    & \tauFzero^{(b)}(X_i) &= \texttt{tau\_f0\_post}[b, i], \\
    \muFone^{(b)}(X_i) &= \texttt{mu\_f1\_post}[b, i],
    & \tauFone^{(b)}(X_i) &= \texttt{tau\_f1\_post}[b, i],
  \end{aligned}
  \label{eq:zi-posterior-storage}
\end{equation}
together with the dispersion draws $\kappa^{(b)}$ (NB / ZINB only).
In the non-ZI sub-model only the first row is stored. Downstream
inference is a post-processing exercise on these matrices and
produces three (count) or four (ZI) classes of summaries.

\paragraph{Per-iteration potential outcomes.}
For each draw $b$ and each unit $i$, the closed-form mapping from
stored forest values to response-scale potential-outcome means is,
for $z \in \{0, 1\}$,
\begin{equation}
  \boxed{\;
    \begin{aligned}
      \log \lambda_z^{(b)}(X_i)
        &\;=\; \muF^{(b)}(X_i) + z \cdot \tauF^{(b)}(X_i), \\[2pt]
      \etazi[z]^{(b)}(X_i)
        &\;=\; \bigl[\muFzero^{(b)}(X_i) + z \cdot \tauFzero^{(b)}(X_i)\bigr]
              - \bigl[\muFone^{(b)}(X_i)  + z \cdot \tauFone^{(b)}(X_i)\bigr], \\[2pt]
      \muz^{(b)}(X_i)
        &\;=\; \bigl(1 - \sigma(\etazi[z]^{(b)}(X_i))\bigr) \,
              \exp\!\bigl(\log \lambda_z^{(b)}(X_i)\bigr).
    \end{aligned}
  \;}
  \label{eq:zi-potoutcomes}
\end{equation}
In the non-ZI sub-model $\sigma(\etazi) \equiv 0$ and
\eqref{eq:zi-potoutcomes} collapses to
$\muz^{(b)}(X_i) = \exp(\muF^{(b)}(X_i) + z \tauF^{(b)}(X_i))$. The
pair $(\muzero^{(b)}, \muone^{(b)})$ collects the response-scale
potential-outcome means at every unit under every draw, and all
reported quantities below are deterministic functions of these two
matrices -- no additional sampling is required: posterior uncertainty
in the response-scale CATE arises from joint uncertainty in the
forests and is automatically propagated by the backfitting Gibbs
sampler.

\paragraph{Response-scale CATE.}
The primary estimand is \eqref{eq:cate-ate}. The $b$th posterior draw
is
\begin{equation}
  \widehat{\tauresp,_i}^{(b)}
  \;=\; \muone^{(b)}(X_i) - \muzero^{(b)}(X_i),
  \label{eq:zi-tauresp-draw}
\end{equation}
with the count-only form $\widehat{\tauresp,_i}^{(b)} =
\exp(\muF^{(b)}(X_i))(\exp(\tauF^{(b)}(X_i)) - 1)$ recovered in the
non-ZI sub-model. The per-unit posterior mean and $95\%$ equal-tailed
credible interval are
$\overline{\tauresp,_i} = B^{-1} \sum_b \widehat{\tauresp,_i}^{(b)}$
and $[Q_{0.025}, Q_{0.975}]$ across draws, with $Q_q$ the empirical
$q$-quantile. The log-rate CATE
$\widehat{\taulog,_i}^{(b)} = \tauF^{(b)}(X_i)$ is reported alongside
for users whose applied questions are best phrased as relative rates.

\paragraph{Sample-average treatment effect.}
The ATE \eqref{eq:cate-ate} is the in-sample mean of the
response-scale CATE,
\begin{equation}
  \widehat{\mathrm{ATE}}^{(b)}
  \;=\; \frac{1}{n} \sum_{i=1}^{n} \widehat{\tauresp,_i}^{(b)},
  \label{eq:ate-draw}
\end{equation}
giving a $B$-vector of posterior draws. The posterior mean and
$95\%$ equal-tailed interval of this vector give the point estimate
and uncertainty of the sample ATE. The average is over the empirical
distribution of $X$, so \eqref{eq:ate-draw} estimates the
\emph{sample-average} treatment effect rather than a
population-average one; \citet{hahn2020bayesian} discuss why this is
the natural target for BART-family methods.

\paragraph{Channel decomposition (ZI only).}
A central affordance of the six-forest construction over a
single-channel causal forest is that $\tauresp$ can be decomposed
into a count-rate channel and a structural-zero channel. The per-unit
posterior draws of the two channel components introduced in
\eqref{eq:channels} are
\begin{equation}
  \widehat{\Deltalam,_i}^{(b)}
    \;=\; \exp\!\bigl(\log\lambda_1^{(b)}(X_i)\bigr)
        - \exp\!\bigl(\log\lambda_0^{(b)}(X_i)\bigr),
  \qquad
  \widehat{\Deltapzi,_i}^{(b)}
    \;=\; \sigma\!\bigl(\etazi[1]^{(b)}(X_i)\bigr)
        - \sigma\!\bigl(\etazi[0]^{(b)}(X_i)\bigr),
  \label{eq:zi-channels-draws}
\end{equation}
both deterministic functions of the stored forest values. Either
channel can be summarised by its posterior mean and $95\%$ credible
interval per unit, or averaged to produce a posterior of the
channel-specific sample average. The two channels are not constrained
to share a sign: for some units the treatment can increase the count
rate \emph{and} decrease the structural-zero probability (both
contributing positively to $\tauresp$); for other units the channels
can offset, producing a small response-scale effect masking large
channel-specific shifts.

\paragraph{Counterfactual prediction at new $X$ and diagnostics.}
For prediction at a new covariate matrix $X^\star$ and treatment
vector $Z^\star$, the user passes \verb|return_trees = TRUE| to
retain the serialised tree traces; \texttt{get\_forest\_fit} then
evaluates each stored tree at $X^\star$, after which the closed-form
mappings above apply unchanged. The default
(\verb|return_trees = FALSE|) drops the tree traces after the
in-sample fit to keep memory consumption modest. All per-unit
summaries are reported on the original (unsorted) input order; the
internal sort that places zeros first (used to amortise the
latent-variable scan) is undone at the storage step
\eqref{eq:zi-posterior-storage}. The NB / ZINB acceptance rate
$\texttt{kappa\_acpt}$ and chain-level traces of $\kappa^{(b)}$ serve
as the main convergence diagnostics; we additionally monitor
effective sample sizes of $\widehat{\mathrm{ATE}}^{(b)}$ and of a
small panel of $\widehat{\tauresp,_i}^{(b)}$ at sentinel units.

% =======================================================================
\section{Research methodology}\label{sec:methodology}

This section pre-registers the simulation study that will evaluate
CountBCF against the natural Gaussian-likelihood baseline. The
specification is held fixed before any cells are run; results will be
reported in Section~\ref{sec:empirical} without revision to the
design below. The motivation for pre-registration is standard
\citep[\S~3.2]{hahn2020bayesian}: heterogeneous-effect comparisons
are sensitive to which DGPs, factor sweeps, and metrics one chooses
to highlight, and reporting those choices \emph{after} seeing results
risks confirming the preferred method.

\subsection{Estimand}\label{sec:method-estimand}

We target the response-scale conditional and average treatment
effects defined in Section~\ref{sec:notation}: for unit $i$, the
potential-outcome means
$\muz(x) = (1 - \pzi(x)) \exp(\log\lambda_z(x))$ of
\eqref{eq:potmeans} give rise to
\begin{equation}
  \tauresp(x) \;=\; \muone(x) - \muzero(x),
  \qquad
  \mathrm{ATE} \;=\; \E_{X}\!\bigl[\tauresp(X)\bigr],
  \label{eq:method-estimand}
\end{equation}
on the same units as $Y$. Under zero-inflation \eqref{eq:method-estimand}
mixes a change in the count rate and a change in the structural-zero
probability, and a method that changes only one of the two will
systematically under-estimate the effect on at least some part of the
covariate space; the ZI DGPs of Section~\ref{sec:method-dgp}
engineer this asymmetry deliberately. All methods, regardless of
their native scale, are scored against this rate-scale truth. For
CountBCF the per-draw CATE is obtained from
\eqref{eq:zi-potoutcomes}--\eqref{eq:zi-tauresp-draw}; for the
Gaussian BCF baseline it is read off directly from the moderating
forest.

\subsection{Data-generating processes}\label{sec:method-dgp}

The simulation grid is a $2 \times 4$ design over (i) the
\emph{functional form} of the prognostic and moderating surfaces
(linear or nonlinear) and (ii) the \emph{count distribution}
(Poisson, NB with $\kappa = 2$, ZIP, or ZINB with $\kappa = 2$). All
eight DGPs share the same covariate distribution and propensity
model, so cross-DGP comparisons isolate the effects of nonlinearity,
overdispersion, and zero-inflation. Each DGP features strong,
sign-changing, multi-covariate heterogeneity; the ZI DGPs additionally
contain non-trivial heterogeneity on the structural-zero channel, so
the per-unit CATE estimation task is substantively challenging in
every cell.

\paragraph{Covariates and treatment.}
We draw $n$ iid covariate vectors $X_i \in \mathbb{R}^p$ with the
first $p - 2$ columns iid standard Gaussian and the last two binary,
\begin{equation}
  X_{i, j} \stackrel{\mathrm{iid}}{\sim} \N(0, 1) \text{ for } j = 1, \dots, p - 2,
  \quad
  X_{i,\, p-1} \sim \Ber(0.6),
  \quad
  X_{i,\, p} \sim \Ber(0.4),
  \label{eq:method-covariates}
\end{equation}
so every design contains two binary columns regardless of $p$. Only
the first five covariates carry signal; columns $X_{\cdot, 6}, \dots,
X_{\cdot, p - 2}$ are irrelevant noise included to stress-test
variable selection as $p$ grows. The propensity model is fixed
across DGPs at
\begin{equation}
  \pi(x) \;=\; \sigma(0.5\, x_1 - 0.4\, x_2),
  \qquad
  Z_i \mid X_i \sim \Ber(\pi(X_i)),
  \label{eq:method-propensity}
\end{equation}
with $\sigma(\cdot)$ the logistic CDF. The estimated propensity
$\pihat$ passed to the methods is set to the truth, $\pihat(X_i) =
\pi(X_i)$, so the comparison isolates the outcome model rather than
penalising methods for propensity mis-estimation.

\paragraph{Count component.}
For all eight DGPs the conditional log-mean of the count component
(non-ZI cells) or the not-structural-zero component (ZI cells) has
the BCF-style form
\begin{equation}
  \log \lambda_i \;=\; \mu_c(X_i) \;+\; Z_i\,\bigl[\tau_0 + \tau_h(X_i)\bigr],
  \label{eq:method-loglambda}
\end{equation}
with the prognostic surface $\mu_c$ and the heterogeneous part
$\tau_h$ taking the form in the upper half of
Table~\ref{tab:method-dgp}. The scalar $\tau_0$ is a per-cell
intercept calibrated so that the realised response-scale ATE matches
the cell's target value. For non-ZI cells, we solve
\begin{equation}
  \E_{X}\!\bigl[\exp(\mu_c(X))\,\bigl(\exp(\tau_0 + \tau_h(X)) - 1\bigr)\bigr]
  \;=\; \mathrm{ATE}_{\mathrm{target}}
  \label{eq:method-tau0-calib}
\end{equation}
via \texttt{R}'s \texttt{uniroot()} on a fresh held-out sample of
$5{,}000$ covariate vectors. For ZI cells, the calibration accounts
for the $(1 - \pzi)$ thinning of the count component; only $\tau_0$
is calibrated, the ZI shift $\Detazi$ is held fixed.

\paragraph{Zero-inflation component (ZI cells only).}
The structural-zero probability follows $\pzi(x) = \sigma(\etazi(x) +
z \Detazi(x))$, with the control logit $\etazi(x)$ and treatment
shift $\Detazi(x)$ in the lower half of Table~\ref{tab:method-dgp}.
The control logit is calibrated so the baseline structural-zero
probability is approximately $25$--$30\%$ across the covariate
distribution -- large enough to make the ZI channel matter but small
enough not to swamp the count signal. The treatment shift
$\Detazi(x)$ is sign-changing across units: for some units the
treatment increases the structural-zero probability while for others
it decreases it, with the sign and magnitude depending on $X$. A
single ATE-level summary of the ZI channel would average over the
sign reversal and report a misleading null; the channel-decomposition
plots described below are designed to recover this heterogeneity.

\begin{table}[t]
\centering
\caption{Functional forms for the prognostic surface $\mu_c$ and the
heterogeneous component $\tau_h$ (count channel, top), and for the
control logit $\etazi$ and treatment shift $\Detazi$ on the ZI
channel (bottom). The linear and nonlinear families each cross with
the four count distributions $\{\Poi, \NB, \mathrm{ZIP},
\mathrm{ZINB}\}$ to give eight DGPs; the ZI rows apply only to the
ZIP / ZINB cells.}
\label{tab:method-dgp}
\begin{tabular}{lll}
\toprule
DGP family & linear & nonlinear \\
\midrule
\multicolumn{3}{l}{\emph{Count component} ($\log\lambda_i = \mu_c(X_i) + Z_i[\tau_0 + \tau_h(X_i)]$)} \\
$\mu_c(x)$ &
  $1.0 + 0.5 x_1 - 0.3 x_2 + 0.2 x_3$ &
  $0.8 + 0.6 \sin(x_1) + 0.4 x_2 x_3 + 0.5\,\ind\{x_4 > 0\}$ \\
$\tau_h(x)$ &
  $0.60 x_1 - 0.45 x_2 + 0.30 x_3$ &
  $0.80 \sin(x_1) + 0.50\,\ind\{x_2 > 0\}\, x_3 - 0.40 x_5$ \\
\midrule
\multicolumn{3}{l}{\emph{ZI component} ($\pzi(x) = \sigma(\etazi(x) + z \Detazi(x))$, ZIP / ZINB only)} \\
$\etazi(x)$ &
  $-1.0 + 0.5 x_2$ &
  $-1.0 + 0.7 \lvert x_2 \rvert - 0.5 \exp(-x_3^{2})$ \\
$\Detazi(x)$ &
  $-0.50 x_1 + 0.30 x_3 - 0.40 x_4$ &
  $-0.60 x_1 x_4 + 0.40 x_2\,\ind\{x_5 > 0\}$ \\
\bottomrule
\end{tabular}
\end{table}

\paragraph{Sampling distribution.}
For each unit $i$, the response is drawn from one of four
distributions. Under the non-ZI cells, $Y_i \sim \Poi(\lambda_i)$ or
$\NB(\kappa, \lambda_i)$ with $\kappa = 2$. Under the ZI cells, a
Bernoulli structural-zero indicator $S_i \sim \Ber(\pzi[Z_i](X_i))$
is drawn first; if $S_i = 1$ then $Y_i = 0$ deterministically;
otherwise
$Y_i \mid S_i = 0 \sim \Poi(\lambda_i)$ or $\NB(\kappa, \lambda_i)$
with $\kappa = 2$. Combined with Table~\ref{tab:method-dgp} this
yields the eight CATE-focused DGPs
\verb|linear_{poisson,nb,zip,zinb}_cate| and
\verb|nonlinear_{poisson,nb,zip,zinb}_cate|.

\subsection{Experimental factors}\label{sec:method-factors}

We sweep three experimental factors -- sample size $N$, covariate
count $p$, and target ATE -- one-at-a-time (OAT) around a reference
cell. Reference values are $N_{\mathrm{ref}} = 250$,
$p_{\mathrm{ref}} = 5$, $\mathrm{ATE}_{\mathrm{ref}} = 1.25$.
Off-reference levels are in Table~\ref{tab:method-factors}; together
with the reference, this gives $1 + 2 + 2 + 2 = 7$ cells per
$(\text{method}, \text{DGP})$ pair. For each cell all eight DGPs of
Section~\ref{sec:method-dgp} are run. Within a cell $\times$ DGP we
draw $N_{\mathrm{sim}} = 100$ independent Monte Carlo replicates;
the seed of replicate $s$ is exactly $s$ -- \verb|set.seed(sim_id)| --
so any single replicate is reproducible from its
$(\text{dgp}, N, p, \text{ATE}, \text{sim\_id})$ tuple alone.

\begin{table}[t]
\centering
\caption{One-at-a-time factor sweep. ``Held fixed'' columns list the
factors that take their reference value within the row's sweep.}
\label{tab:method-factors}
\begin{tabular}{lll}
\toprule
Factor & Off-reference levels & Held fixed \\
\midrule
Sample size $N$         & $\{100,\, 500\}$       & $p = 5$,\; ATE $= 1.25$ \\
Covariate count $p$     & $\{50,\, 250\}$        & $N = 250$,\; ATE $= 1.25$ \\
Target ATE              & $\{0.5,\, 2.5\}$       & $N = 250$,\; $p = 5$ \\
\bottomrule
\end{tabular}
\end{table}

\subsection{Methods compared}\label{sec:method-methods}

Both methods produce a posterior over the response-scale ATE and
per-unit posterior means and credible intervals for the
response-scale CATE; both are Bayesian, both share the BCF prior
structure, and they differ only in the outcome likelihood.

\begin{enumerate}
  \item \textbf{CountBCF (proposed).} Fit via \verb|countbcf()| with
        \verb|count_model| set to one of \verb|"poisson"|,
        \verb|"nb"|, \verb|"zipoisson"|, or \verb|"zinb"| depending
        on the DGP. All other hyperparameters take the defaults
        catalogued in Table~\ref{tab:default-hyperparams}. The
        per-unit CATE is recovered from the closed-form mapping
        \eqref{eq:zi-potoutcomes}--\eqref{eq:zi-tauresp-draw}; under
        ZIP / ZINB the same coefficient matrices are also
        post-processed into the channel decomposition
        $(\widehat\Deltalam, \widehat\Deltapzi)$ of
        \eqref{eq:zi-channels-draws}.
  \item \textbf{BCF (Gaussian baseline).} The
        \citet{hahn2020bayesian} BCF applied to the count outcome
        \emph{as if it were Gaussian}, fit via \verb|bcf_binary()|.
        The treatment-effect posterior is read off the moderating
        forest with \verb|get_forest_fit()|. This is the natural
        ``what if you ignored the count structure'' baseline: it
        retains the BCF $(\mu, \tau)$ decomposition and the
        propensity-as-covariate fix but does not respect the
        discreteness of $Y$, the heteroscedasticity induced by the
        count likelihood, or the structural-zero process; in the ZI
        case it produces only a single response-scale CATE with no
        channel decomposition.
\end{enumerate}

Both methods use the same MCMC budget: $\mathtt{nburn} =
\mathtt{nsim} = 1000$, $\mathtt{nthin} = 1$. The methodological
symmetry is deliberate -- the only structural difference is the
outcome likelihood -- so any performance gap observed in
Section~\ref{sec:empirical} can be attributed to the likelihood, not
to incidental choices of prior, MCMC budget, or post-processing.

\paragraph{CountBART considered and dropped.}
An $S$-learner built on top of \verb|count_bart()| -- augmenting the
design matrix with $Z$ and $\pihat$, fitting a single log-linear BART
model on $(X, Z, \pihat)$ (three-forest in the ZI case), and reading
the CATE off counterfactual evaluations -- was originally part of the
benchmark. It consistently exhausts the available Colab RAM and
crashes on the larger ablation cells ($p \in \{50, 250\}$); the
three-forest backbone of the ZI model amplifies the memory pressure
relative to the non-ZI case. Until the memory footprint of the
upstream \texttt{count\_bart} is reduced, the CountBART $S$-learner is
not feasible to evaluate against CountBCF on the same hardware, and
we exclude it rather than report it on the subset of cells where it
runs.

\subsection{Metrics}\label{sec:method-metrics}

For each row of every per-cell CSV -- one row per
$(\text{cell}, \text{DGP}, \text{sim}, \text{method})$ tuple -- we
record three families of metrics plus light bookkeeping. All metrics
that involve the truth are computed on the \emph{response} scale,
regardless of which scale the method natively reports. Aggregating
across the $N_{\mathrm{sim}} = 100$ replicates within a cell gives
Monte Carlo means and standard errors. The CATE-level metrics are
the primary site of method differentiation in the $N = 100$ regime.

\paragraph{Outcome-fit, ATE, and CATE.}
Two diagnostics measure how well the fitted mean approximates the
true conditional mean on the observed arm:
$\mathtt{rmse\_yhat} = \sqrt{\overline{(\widehat y_i - \mu_{Z_i}(X_i))^2}}$
and $\mathtt{bias\_yhat} = \overline{\widehat y_i - \mu_{Z_i}(X_i)}$,
with $\widehat y_i$ the posterior-mean fitted value under the
observed treatment and $\mu_{Z_i}(X_i) = (1 - \pzi[Z_i](X_i))
\exp(\log\lambda_{Z_i}(X_i))$ the response-scale truth. ATE-level
inference is summarised by \verb|ate_mean|, \verb|ate_q025|,
\verb|ate_q975| (the $95\%$ credible-interval bounds),
\verb|ate_coverage| (the indicator that the interval contains the
realised true ATE), and \verb|true_ate|
($n^{-1}\sum_i \tauresp(X_i)$). CATE-level performance is summarised
by the seven metrics of Table~\ref{tab:method-metrics-cate}.

\begin{table}[t]
\centering
\caption{CATE-level metrics, evaluated on the response scale against
the realised truth $\tau(X_i) = \muone(X_i) - \muzero(X_i)$.}
\label{tab:method-metrics-cate}
\begin{tabular}{ll}
\toprule
Metric & Definition \\
\midrule
\verb|pehe| &
  $\sqrt{\,\overline{(\widehat\tau(X_i) - \tau(X_i))^{2}}\,}$ (PEHE) \\
\verb|mae_cate| &
  $\overline{\lvert \widehat\tau(X_i) - \tau(X_i) \rvert}$ \\
\verb|cate_bias| &
  $\overline{\widehat\tau(X_i) - \tau(X_i)}$ \\
\verb|cate_cor| &
  Pearson correlation between $\widehat\tau$ and $\tau$ \\
\verb|cate_cov95|, \verb|cate_cov50| &
  Fraction of units whose 95\% / 50\% unit-level CI contains $\tau(X_i)$ \\
\verb|cate_ci_width95| &
  Mean width of the unit-level 95\% credible intervals \\
\verb|cate_sd_true| &
  $\mathrm{sd}(\tau(X_i))$ across the in-sample units \\
\bottomrule
\end{tabular}
\end{table}

PEHE \citep{hill2011bayesian} is the headline metric, but in the
$N_{\mathrm{sim}} = 100$ regime the credible-interval metrics --
particularly \verb|cate_cov95| and \verb|cate_ci_width95| -- are what
most clearly expose miscalibration of the Gaussian baseline.

\paragraph{Zero-inflation regime indicators (ZI cells only).}
Two ZI-specific columns record the regime each Monte Carlo replicate
sits in:
$\mathtt{pct\_zero} = n^{-1}\sum_i \ind\{Y_i = 0\}$
(the realised total-zero rate, observable from $Y$) and
$\mathtt{pct\_struct\_zero} = n^{-1}\sum_i \ind\{S_i = 1\}$
(the realised structural-zero rate, available only via the latent
indicator drawn by the DGP). Together with their ratio, these
characterise how much of the observed-zero mass is structural versus
count-process at the cell's calibration point, and make it possible
to distinguish performance differences driven by the ZI channel from
those driven by overall zero-rate effects.

\paragraph{Bookkeeping.}
Two columns are kept for transparency: \verb|elapsed_sec| (the
wall-clock fit time, in seconds) and \verb|error| (\verb|NA| on
success; the caught \verb|R| error message otherwise, so failure
modes are recorded rather than silently dropped).

\subsection{Outputs and reproducibility}\label{sec:method-outputs}

Each $(\text{method}, \text{DGP}, \text{cell})$ tuple writes its own
long-format CSV and a timestamped progress log. CSV filenames follow
the pattern
\verb|sim__<method>__<dgp>__N<N>_P<P>_ATE<ATE x100>.csv|, with
companion \verb|.log| files. There are $2$ methods $\times$ $8$ CATE
DGPs $\times$ $7$ cells $= 112$ CSVs across the count-only and ZI
sub-studies combined, each holding $100$ rows (one per Monte Carlo
replicate). Rows are streamed incrementally so that a crash on
replicate $s$ does not lose the work of replicates $1, \dots, s - 1$.
Aggregation is a one-liner: \verb|rbind| of the CSVs in \verb|R| or
\verb|cat *.csv| in a shell.

\paragraph{Output panels.}
Section~\ref{sec:empirical} will present three families of panels.
\emph{Outcome-fit / ATE / CATE} panels show the standard Monte Carlo
summaries of the metrics in Section~\ref{sec:method-metrics} along
the three OAT sweeps. \emph{Coverage diagnostics} report empirical
$95\%$ interval coverage of both \verb|ate_coverage| and
\verb|cate_cov95| so that miscalibration is visible on both the
marginal and conditional scales. \emph{Channel decomposition} -- the
panel template specific to the ZI sub-study -- plots, for CountBCF
only, the per-unit posterior means of $\Deltalam$ and $\Deltapzi$
against each other and against the combined response-scale CATE
$\tauresp$ on a per-DGP basis at the reference cell. Channels that
share sign across most of the covariate space appear as a clustered
cloud along the diagonal; sign-changing DGPs produce a cross-shaped
scatter; either pattern is a visual confirmation that the six-forest
decomposition recovers heterogeneity that single-channel methods
cannot expose.

\paragraph{Reproducibility.}
The runners under \verb|tests/runners_count/| and
\verb|tests/runners_zi/| are self-contained single \verb|.R| files:
each depends only on \verb|library(countbart)| and base \verb|R| --
no \verb|source()| of any repository file -- so any single runner is
portable to a Google Colab session without setup. With $15$ CPUs the
runners execute in waves of $15$. The runners themselves are emitted
by \verb|tests/_generate_runners.R|, which holds the cell grid, the
DGP function bodies, and the method wrappers as templated strings;
editing it and re-running regenerates the runners in lock-step. Each
replicate seed equals its \verb|sim_id|, so any individual replicate
is reproducible from its
$(\text{dgp}, N, p, \text{ATE}, \text{sim\_id})$ tuple alone. We do
not anonymise or post-process the per-cell CSVs prior to plotting:
the raw output is the audit trail.

% =======================================================================
\section{Simulation Studies}\label{sec:empirical}

This section reports the outcome of the pre-registered protocol of
Section~\ref{sec:methodology}, with no revision to the design. Two
sub-studies are run in parallel: a \emph{count-only} sub-study over
the four CATE-focused Poisson and Negative-Binomial DGPs of
Table~\ref{tab:method-dgp}, and a \emph{zero-inflated} sub-study over
the corresponding ZIP and ZINB cells. The two sub-studies isolate
complementary phenomena. The count-only sub-study measures the cost
of treating an integer outcome as Gaussian when no structural-zero
channel is present; the zero-inflated sub-study additionally stresses
the methods on a second treatment channel ($\Detazi$) that the
Gaussian baseline cannot represent at all. Because both sub-studies
share the covariate distribution, propensity model
\eqref{eq:method-propensity}, factor grid
(Table~\ref{tab:method-factors}), and Monte Carlo budget
($N_{\mathrm{sim}} = 100$ replicates per cell), every contrast
across the two sub-studies can be attributed to the presence or
absence of zero-inflation in the DGP, not to incidental design
choices.

The remainder of the section is organised in three parts. The
reference-cell head-to-head (\S~\ref{sec:emp-ref}) fixes the design
at $(N, p, \mathrm{ATE}) = (250, 5, 1.25)$ and reports both methods
on all six metrics across all eight DGPs. The ablation studies
(\S~\ref{sec:emp-abl}) sweep each experimental factor one at a time
around the reference cell and present, for each sweep, the three
CATE-level metrics most diagnostic of the comparison: the per-unit
$95\%$ credible-interval coverage $\mathtt{cate\_cov95}$, the
heterogeneous-effect error PEHE, and the average per-unit $95\%$
interval width $\mathtt{cate\_ci\_width95}$. The cross-cutting
interpretation (\S~\ref{sec:emp-interp}) and a list of caveats
(\S~\ref{sec:emp-caveats}) close the section. Aggregate ATE-level
quantities ($\mathtt{rmse\_ate}$, $\mathtt{ate\_coverage}$, ATE CI
width) and RMSE / ATE coverage curves along the same sweeps are
omitted from the figures to keep the visual presentation focused on
the heterogeneity story; the underlying per-cell numbers are
provided in \texttt{tests/analysis/summary\_count.csv} and
\texttt{tests/analysis/summary\_zi.csv}, and the headline numbers
are quoted in the text.

\subsection{Reference-cell head-to-head}\label{sec:emp-ref}

We first fix the design at the reference cell
$(N, p, \mathrm{ATE}) = (250, 5, 1.25)$ and compare the two methods
at the centre of the factor grid. Tables~\ref{tab:sim_count_ref}
and~\ref{tab:sim_zi_ref} report all six metrics for the four
count-only and four zero-inflated DGPs respectively. RMSE(ATE) and
ATE coverage are scalar Monte Carlo summaries; the remaining columns
are reported as cross-replicate mean $\pm$ one standard deviation
across the $100$ replicates per cell.

\begin{table}[h!]\footnotesize
\caption{Reference-cell ($N = 250$, $p = 5$, ATE $= 1.25$) Monte-Carlo summary on the count-only DGPs (Poisson, NB). RMSE ATE and ATE coverage are scalar Monte-Carlo summaries; the remaining columns report the cross-replicate mean $\pm$ one standard deviation. ATE CI Width is the average width of the 95\% ATE posterior credible interval; PEHE is $\sqrt{\overline{(\hat\tau_i - \tau_i)^2}}$; CATE Cov.~95 is the per-unit 95\% credible-interval coverage of $\tau(X_i)$; CATE CI Width95 is the mean width of those unit-level intervals.}\label{tab:sim_count_ref}
\centering
\begin{tabular*}{\linewidth}{ @{\extracolsep{\fill}} ll cc cccc @{}}
\toprule
 &  & \multicolumn{2}{c}{ATE} & \multicolumn{4}{c}{CATE} \\
\cmidrule(lr){3-4} \cmidrule(lr){5-8}
DGP & Method & RMSE & Cov.\,95 & Avg.\ CI Width & PEHE & Cov.\,95 & Avg.\ CI Width95 \\
\midrule
\multirow{2}{*}{Linear / Poisson} & BCF (Gaussian) & 1.15 & 0.55 & 1.23 $\pm$ 0.26 & 6.23 $\pm$ 4.44 & 0.81 $\pm$ 0.08 & 4.20 $\pm$ 1.12 \\
 & CountBCF & 1.08 & 0.65 & 2.28 $\pm$ 0.82 & 4.77 $\pm$ 2.46 & 0.92 $\pm$ 0.07 & 7.83 $\pm$ 2.21 \\
\midrule
\multirow{2}{*}{Nonlinear / Poisson} & BCF (Gaussian) & 1.09 & 0.64 & 1.40 $\pm$ 0.18 & 5.92 $\pm$ 4.80 & 0.81 $\pm$ 0.07 & 5.32 $\pm$ 0.98 \\
 & CountBCF & 0.76 & 0.80 & 1.88 $\pm$ 0.44 & 3.48 $\pm$ 1.66 & 0.87 $\pm$ 0.09 & 6.74 $\pm$ 1.75 \\
\midrule
\multirow{2}{*}{Linear / NB} & BCF (Gaussian) & 1.21 & 0.60 & 1.79 $\pm$ 0.36 & 9.16 $\pm$ 4.90 & 0.82 $\pm$ 0.06 & 5.54 $\pm$ 1.63 \\
 & CountBCF & 1.25 & 0.65 & 2.95 $\pm$ 1.18 & 5.95 $\pm$ 2.83 & 0.90 $\pm$ 0.07 & 10.08 $\pm$ 3.09 \\
\midrule
\multirow{2}{*}{Nonlinear / NB} & BCF (Gaussian) & 1.30 & 0.68 & 2.23 $\pm$ 0.57 & 7.81 $\pm$ 5.95 & 0.83 $\pm$ 0.08 & 7.34 $\pm$ 2.07 \\
 & CountBCF & 0.90 & 0.80 & 2.63 $\pm$ 0.65 & 4.49 $\pm$ 1.81 & 0.88 $\pm$ 0.07 & 9.61 $\pm$ 2.95 \\
\bottomrule
\end{tabular*}
\end{table}

\begin{table}[h!]\footnotesize
\caption{Reference-cell ($N = 250$, $p = 5$, ATE $= 1.25$) Monte-Carlo summary on the zero-inflated DGPs (ZIP, ZINB). Conventions match Table~\ref{tab:sim_count_ref}.}\label{tab:sim_zi_ref}
\centering
\begin{tabular*}{\linewidth}{ @{\extracolsep{\fill}} ll cc cccc @{}}
\toprule
 &  & \multicolumn{2}{c}{ATE} & \multicolumn{4}{c}{CATE} \\
\cmidrule(lr){3-4} \cmidrule(lr){5-8}
DGP & Method & RMSE & Cov.\,95 & Avg.\ CI Width & PEHE & Cov.\,95 & Avg.\ CI Width95 \\
\midrule
\multirow{2}{*}{Linear / ZIP} & BCF (Gaussian) & 0.97 & 0.62 & 1.25 $\pm$ 0.23 & 5.81 $\pm$ 4.03 & 0.83 $\pm$ 0.07 & 4.06 $\pm$ 1.20 \\
 & CountBCF & 0.73 & 0.76 & 1.88 $\pm$ 0.56 & 3.96 $\pm$ 1.89 & 0.93 $\pm$ 0.05 & 6.67 $\pm$ 1.67 \\
\midrule
\multirow{2}{*}{Nonlinear / ZIP} & BCF (Gaussian) & 1.03 & 0.74 & 1.61 $\pm$ 0.25 & 5.52 $\pm$ 4.41 & 0.83 $\pm$ 0.07 & 5.81 $\pm$ 1.22 \\
 & CountBCF & 0.55 & 0.83 & 1.69 $\pm$ 0.35 & 3.06 $\pm$ 1.47 & 0.93 $\pm$ 0.04 & 6.92 $\pm$ 1.47 \\
\midrule
\multirow{2}{*}{Linear / ZINB} & BCF (Gaussian) & 1.71 & 0.54 & 1.62 $\pm$ 0.35 & 8.45 $\pm$ 9.24 & 0.83 $\pm$ 0.08 & 4.82 $\pm$ 1.78 \\
 & CountBCF & 0.95 & 0.70 & 2.39 $\pm$ 0.72 & 5.17 $\pm$ 2.50 & 0.94 $\pm$ 0.06 & 8.35 $\pm$ 2.33 \\
\midrule
\multirow{2}{*}{Nonlinear / ZINB} & BCF (Gaussian) & 1.35 & 0.74 & 2.12 $\pm$ 0.38 & 7.58 $\pm$ 6.24 & 0.84 $\pm$ 0.08 & 6.77 $\pm$ 1.60 \\
 & CountBCF & 0.67 & 0.91 & 2.36 $\pm$ 0.51 & 3.49 $\pm$ 1.27 & 0.93 $\pm$ 0.04 & 9.02 $\pm$ 2.38 \\
\bottomrule
\end{tabular*}
\end{table}

\paragraph{Aggregate gap.} Averaged over the four DGPs within each
sub-study, CountBCF improves on the Gaussian BCF baseline on every
aggregate metric in both sub-studies. On the count-only DGPs the
average ATE root-mean-squared error falls from $1.19$ to $1.00$
($-16\%$) and the average PEHE from $7.28$ to $4.67$ ($-36\%$); the
corresponding gains on the zero-inflated DGPs are roughly twice as
large in relative terms, $1.26 \to 0.73$ on RMSE(ATE) ($-43\%$) and
$6.84 \to 3.92$ on PEHE ($-43\%$). The gap is widest precisely where
the count likelihood differs most from a Gaussian: the nonlinear
ZINB cell yields the most extreme PEHE ratio
($7.58 \to 3.49$, a $-54\%$ reduction), and the linear-Poisson cell
--- the cell closest to a homoscedastic-Gaussian regime --- the
smallest ($-23\%$). That ordering is direct evidence that the
likelihood, not the BCF prior structure shared by the two methods,
drives the improvement: where the Gaussian working model is least
wrong (linear, Poisson, no zero-inflation), the gap shrinks;
elsewhere it widens.

\paragraph{Coverage and sharpness.} The two methods produce
qualitatively different posterior intervals. Gaussian BCF intervals
are uniformly narrower (count: average ATE CI $1.66$ vs.\ $2.44$,
$-32\%$; ZI: $1.65$ vs.\ $2.08$, $-21\%$), but their empirical
$95\%$ coverage sits at $0.62$ (count) and $0.66$ (ZI) on average
across DGPs --- and falls as low as $0.54$ on the linear ZINB cell,
or $40$ percentage points below nominal. CountBCF widens both
ATE-level and per-unit intervals, paying a sharpness cost of $+46\%$
(count) and $+26\%$ (ZI) on average ATE width and $+53\%$ /\ $+44\%$
on average per-unit CATE width, in exchange for moving empirical
coverage to $0.73$ / $0.80$ on ATE and to $0.89$ / $0.93$ on
per-unit CATE. The latter is essentially at nominal: across the
four zero-inflated DGPs the per-unit CATE coverage of CountBCF lies
in $[0.93, 0.94]$, while the Gaussian baseline lies in
$[0.83, 0.84]$. Under-coverage of a posterior interval is the more
serious failure mode for downstream policy inference --- a $0.83$
interval that one reads as a $0.95$ interval understates the
estimation risk by a factor of nearly three on the tail probability
--- so the wider intervals are a favourable trade.

\paragraph{Per-DGP wins.} Looking at the eight DGPs individually,
CountBCF wins on PEHE in $8/8$ cells (count: $4/4$; ZI: $4/4$) and
is strictly closer to nominal $0.95$ per-unit CATE coverage in
$8/8$. On ATE RMSE the picture is similar ($3/4$ on count, $4/4$ on
ZI); the single tie/loss is on \verb|linear_nb_cate|
($1.21$ vs.\ $1.25$), a difference of order $\widehat\sigma /
\sqrt{100}$ given the per-replicate dispersion and within Monte
Carlo noise. The largest single per-DGP gap is again on the
nonlinear ZINB cell: ATE RMSE $1.35 \to 0.67$ ($-50\%$), PEHE
$7.58 \to 3.49$ ($-54\%$), ATE coverage $0.74 \to 0.91$, per-unit
CATE coverage $0.84 \to 0.93$. This is the DGP that combines all
three pathologies CountBCF was designed for: overdispersion (NB),
nonlinearity in the prognostic and moderating surfaces, and a
structural-zero channel that the Gaussian baseline cannot represent.

\subsection{Ablation studies}\label{sec:emp-abl}

We now sweep each of the three experimental factors of
Table~\ref{tab:method-factors} one at a time around the reference
cell. Each ablation figure below holds the other two factors at
their reference values and overlays the four CATE-focused DGPs of
the relevant sub-study; the dashed grey horizontal in the coverage
panels marks the nominal $0.95$. Error bars denote one
cross-replicate SD across the $100$ Monte Carlo replicates per cell.
The plotted metrics --- per-unit CATE $95\%$ coverage, PEHE, and
average per-unit CATE $95\%$ interval width --- are the three CATE
diagnostics that most clearly expose the calibration of an HTE
posterior; the corresponding ATE-level curves move in the same
direction and are omitted to keep each figure to three panels. Cell
counts in cross-DGP averages reported in the text exclude the single
missing cell (CountBCF on \verb|nonlinear_poisson_cate| at
$N = 250$, $p = 250$); panel (b) of Figure~\ref{fig:abl-count-P}
therefore shows three of the four CountBCF curves at $p = 250$.

\subsubsection*{Sample size sweep ($N \in \{100, 250, 500\}$)}

The sample-size sweep is the cleanest of the three because both
methods should approach a regime where MCMC posterior intervals
shrink and means concentrate as $N$ grows. The numerics confirm
that pattern in different ways for the two methods. On the
count-only sub-study, RMSE(ATE) falls for both methods as expected
(BCF (Gaussian): $-61\%$, CountBCF: $-30\%$ from $N = 100$ to
$N = 500$, averaged over DGPs), but their PEHE responses to extra
data differ qualitatively: CountBCF's PEHE shrinks by $-6\%$, while
Gaussian BCF's is essentially flat ($-4\%$). The mechanism is
visible in the coverage panels (Figure~\ref{fig:abl-count-N}a): the
Gaussian baseline's intervals shrink faster than its bias, so its
empirical ATE coverage drifts $-22\%$ across the sweep, while
CountBCF's intervals shrink in step and its coverage drifts only
$-10\%$. The same pattern is sharper on the zero-inflated
sub-study (Figure~\ref{fig:abl-zi-N}): Gaussian BCF's PEHE
\emph{worsens} by $+10\%$ as $N$ grows, because the additional data
sharpens the wrong working model, while CountBCF's PEHE shrinks by
$-7\%$ and its per-unit CATE coverage rises monotonically toward
$0.96$ on every DGP.

A second feature of the count-only $N$ sweep deserves comment: at
$N = 100$, CountBCF's per-unit CATE coverage is markedly below
nominal on the Poisson DGPs ($0.66$ on linear Poisson, $0.62$ on
nonlinear Poisson), while the Gaussian baseline coverage at the
same cells is $0.87$. This is the only regime in which the
Gaussian baseline reports nominally higher coverage than CountBCF
on this metric, and it is explained by the small-sample behaviour
of the Pólya--Gamma sampler underlying the count likelihood: with
only $\approx 100$ observations the latent-variable augmentation
mixes more slowly, so the posterior on the leaf parameters has not
yet inflated to its asymptotic width. By $N = 250$ CountBCF
recovers ($0.92$ and $0.87$ respectively) and by $N = 500$ it is
above $0.95$. The Gaussian baseline's apparent coverage advantage
at $N = 100$ on the Poisson DGPs is therefore a feature of two
methods that are both miscalibrated at small $N$ on different
sides; it disappears once the sample is large enough to populate
the prior.

\begin{figure}[t]
\centering
\begin{subfigure}[b]{0.32\linewidth}
  \includegraphics[width=\linewidth]{count_N_cate_cov95_mean.png}
  \caption{Per-unit CATE $95\%$ coverage.}
\end{subfigure}\hfill
\begin{subfigure}[b]{0.32\linewidth}
  \includegraphics[width=\linewidth]{count_N_pehe_mean.png}
  \caption{PEHE.}
\end{subfigure}\hfill
\begin{subfigure}[b]{0.32\linewidth}
  \includegraphics[width=\linewidth]{count_N_cate_ci_width95_mean.png}
  \caption{Avg.\ per-unit CATE $95\%$ CI width.}
\end{subfigure}
\caption{Count-only sub-study, sample-size sweep
$N \in \{100, 250, 500\}$ at $p = 5$, $\mathrm{ATE} = 1.25$. Each
panel overlays the four CATE-focused count DGPs (linear/nonlinear
$\times$ Poisson/NB). Dashed line in panel (a): nominal $0.95$
coverage. Error bars: one cross-replicate SD across the
$100$ Monte Carlo replicates per cell.}
\label{fig:abl-count-N}
\end{figure}

\begin{figure}[t]
\centering
\begin{subfigure}[b]{0.32\linewidth}
  \includegraphics[width=\linewidth]{zi_N_cate_cov95_mean.png}
  \caption{Per-unit CATE $95\%$ coverage.}
\end{subfigure}\hfill
\begin{subfigure}[b]{0.32\linewidth}
  \includegraphics[width=\linewidth]{zi_N_pehe_mean.png}
  \caption{PEHE.}
\end{subfigure}\hfill
\begin{subfigure}[b]{0.32\linewidth}
  \includegraphics[width=\linewidth]{zi_N_cate_ci_width95_mean.png}
  \caption{Avg.\ per-unit CATE $95\%$ CI width.}
\end{subfigure}
\caption{Zero-inflated sub-study, sample-size sweep
$N \in \{100, 250, 500\}$ at $p = 5$, $\mathrm{ATE} = 1.25$. Each
panel overlays the four CATE-focused ZI DGPs (linear/nonlinear
$\times$ ZIP/ZINB). Conventions as in
Figure~\ref{fig:abl-count-N}.}
\label{fig:abl-zi-N}
\end{figure}

\subsubsection*{Covariate-count sweep ($p \in \{5, 50, 250\}$)}

The covariate-count sweep is the sweep that most clearly separates
the two methods. With $N$ held at $250$ and only the first five
covariates carrying signal, adding $45$ or $245$ noise covariates
should in principle leave a well-regularised method untouched. On
the count-only sub-study the Gaussian BCF baseline degrades
substantially as nuisance dimensions are added: averaged across
DGPs, RMSE(ATE) rises by $+72\%$ from $p = 5$ to $p = 250$, PEHE by
$+44\%$, and empirical ATE coverage falls by $-30\%$. CountBCF
moves in the opposite direction on RMSE(ATE), \emph{improving} by
$-34\%$ from $p = 5$ to $p = 250$, and worsens on PEHE by only
$+29\%$ (Figure~\ref{fig:abl-count-P}). Per-unit CATE coverage
under CountBCF stays within $[0.83, 0.86]$ across $p$ on the count
DGPs and within $[0.92, 0.94]$ on the ZI DGPs, while the Gaussian
baseline drifts from $\sim 0.82$ at $p = 5$ to $\sim 0.74$ at
$p = 250$ on the count cells and from $\sim 0.83$ to $\sim 0.77$
on the ZI cells.

The asymmetric response to nuisance covariates is, we conjecture,
the combined effect of two features specific to the count
likelihood. First, the Dirichlet-style splitting rule on the log
scale concentrates split mass on the few covariates that explain
residual count structure; when the response is heavy-tailed,
spurious splits on noise covariates make the likelihood
\emph{worse}, so the Metropolis acceptance ratio penalises them
more aggressively than under a homoscedastic Gaussian likelihood
of the same nominal mean. Second, the propensity-score adjustment
shared by the two methods is itself less informative at large $p$,
because the Gaussian baseline must absorb the residual count
heteroscedasticity into the moderating forest, and the wider the
covariate basis the more noise covariates can stand in as proxies
for that residual structure. CountBCF separates these channels by
construction.

\begin{figure}[t]
\centering
\begin{subfigure}[b]{0.32\linewidth}
  \includegraphics[width=\linewidth]{count_P_cate_cov95_mean.png}
  \caption{Per-unit CATE $95\%$ coverage.}
\end{subfigure}\hfill
\begin{subfigure}[b]{0.32\linewidth}
  \includegraphics[width=\linewidth]{count_P_pehe_mean.png}
  \caption{PEHE.}
\end{subfigure}\hfill
\begin{subfigure}[b]{0.32\linewidth}
  \includegraphics[width=\linewidth]{count_P_cate_ci_width95_mean.png}
  \caption{Avg.\ per-unit CATE $95\%$ CI width.}
\end{subfigure}
\caption{Count-only sub-study, covariate-count sweep
$p \in \{5, 50, 250\}$ at $N = 250$, $\mathrm{ATE} = 1.25$. The
single missing cell (CountBCF on \texttt{nonlinear\_poisson\_cate}
at $p = 250$) is omitted from the corresponding CountBCF curve;
the BCF (Gaussian) curve at $p = 250$ for that DGP is shown.
Conventions as in Figure~\ref{fig:abl-count-N}.}
\label{fig:abl-count-P}
\end{figure}

\begin{figure}[t]
\centering
\begin{subfigure}[b]{0.32\linewidth}
  \includegraphics[width=\linewidth]{zi_P_cate_cov95_mean.png}
  \caption{Per-unit CATE $95\%$ coverage.}
\end{subfigure}\hfill
\begin{subfigure}[b]{0.32\linewidth}
  \includegraphics[width=\linewidth]{zi_P_pehe_mean.png}
  \caption{PEHE.}
\end{subfigure}\hfill
\begin{subfigure}[b]{0.32\linewidth}
  \includegraphics[width=\linewidth]{zi_P_cate_ci_width95_mean.png}
  \caption{Avg.\ per-unit CATE $95\%$ CI width.}
\end{subfigure}
\caption{Zero-inflated sub-study, covariate-count sweep
$p \in \{5, 50, 250\}$ at $N = 250$, $\mathrm{ATE} = 1.25$.
Conventions as in Figure~\ref{fig:abl-count-N}; all $56$ ZI cells
completed $100$ replicates.}
\label{fig:abl-zi-P}
\end{figure}

The same story repeats on the zero-inflated sub-study, where all
$56$ cells completed: RMSE(ATE) for the Gaussian baseline rises by
$+60\%$ across $p \in \{5, 250\}$ averaged across DGPs, while
CountBCF \emph{falls} by $-31\%$. PEHE for Gaussian BCF rises by
$+41\%$; CountBCF by $+18\%$. The largest single contrast is on
\verb|linear_zinb_cate|: at $p = 250$ Gaussian BCF reports RMSE(ATE)
$= 3.49$ against CountBCF's $0.52$ --- nearly a seven-fold gap --- and
empirical ATE coverage of $0.34$ against CountBCF's $0.98$. Two
plausible reasons for the unusually good performance of CountBCF
at the largest $p$ on this cell: the Dirichlet-on-splits prior
concentrates on the five informative covariates, and the ZI channel
provides the model with an additional degree of freedom against
which to assign sign-changing treatment effects.

\subsubsection*{Target-ATE sweep ($\mathrm{ATE} \in \{0.5, 1.25, 2.5\}$)}

The target-ATE sweep stresses the magnitude of the treatment effect
while holding the functional form of $\tau_h$ fixed; recall that
$\tau_0$ is calibrated per cell so that the realised response-scale
ATE matches the target. As the calibrated treatment effect grows
from $0.5$ to $2.5$, both methods' absolute errors increase, as one
would expect from the multiplicative scaling of the count outcome:
on the count-only sub-study, average BCF RMSE rises by $+45\%$,
CountBCF by $+63\%$; BCF PEHE rises by $+54\%$, CountBCF by $+54\%$
(Figure~\ref{fig:abl-count-ATE}). On the zero-inflated sub-study,
both methods show steeper growth: BCF RMSE $+77\%$, CountBCF
$+93\%$; BCF PEHE $+79\%$, CountBCF $+64\%$
(Figure~\ref{fig:abl-zi-ATE}).

The headline contrast is that absolute errors grow proportionally
on both methods, but CountBCF stays below BCF (Gaussian) on RMSE
and PEHE at every level of the swept target ATE, and the
per-unit CATE coverage of CountBCF is the only one of the four
trajectories that stays near nominal as the effect size grows:
across the swept range CountBCF holds at $0.87$--$0.95$ on the
count DGPs and at $0.91$--$0.96$ on the ZI DGPs, while the Gaussian
baseline drops from $\sim 0.85$ at $\mathrm{ATE} = 0.5$ to
$\sim 0.78$ at $\mathrm{ATE} = 2.5$ on both sub-studies. CI widths
in both methods grow proportionally to the calibrated ATE: this
is a property of the response-scale CATE, which inherits the
multiplicative scaling of the count rate, not of either posterior;
it is visible in panel (c) of both ATE-sweep figures as nearly
linear growth on both methods.

\begin{figure}[t]
\centering
\begin{subfigure}[b]{0.32\linewidth}
  \includegraphics[width=\linewidth]{count_ATE_cate_cov95_mean.png}
  \caption{Per-unit CATE $95\%$ coverage.}
\end{subfigure}\hfill
\begin{subfigure}[b]{0.32\linewidth}
  \includegraphics[width=\linewidth]{count_ATE_pehe_mean.png}
  \caption{PEHE.}
\end{subfigure}\hfill
\begin{subfigure}[b]{0.32\linewidth}
  \includegraphics[width=\linewidth]{count_ATE_cate_ci_width95_mean.png}
  \caption{Avg.\ per-unit CATE $95\%$ CI width.}
\end{subfigure}
\caption{Count-only sub-study, target-ATE sweep
$\mathrm{ATE} \in \{0.5, 1.25, 2.5\}$ at $N = 250$, $p = 5$. The
per-cell intercept $\tau_0$ is calibrated so the realised
response-scale ATE matches the target. Conventions as in
Figure~\ref{fig:abl-count-N}.}
\label{fig:abl-count-ATE}
\end{figure}

\begin{figure}[t]
\centering
\begin{subfigure}[b]{0.32\linewidth}
  \includegraphics[width=\linewidth]{zi_ATE_cate_cov95_mean.png}
  \caption{Per-unit CATE $95\%$ coverage.}
\end{subfigure}\hfill
\begin{subfigure}[b]{0.32\linewidth}
  \includegraphics[width=\linewidth]{zi_ATE_pehe_mean.png}
  \caption{PEHE.}
\end{subfigure}\hfill
\begin{subfigure}[b]{0.32\linewidth}
  \includegraphics[width=\linewidth]{zi_ATE_cate_ci_width95_mean.png}
  \caption{Avg.\ per-unit CATE $95\%$ CI width.}
\end{subfigure}
\caption{Zero-inflated sub-study, target-ATE sweep
$\mathrm{ATE} \in \{0.5, 1.25, 2.5\}$ at $N = 250$, $p = 5$.
Conventions as in Figure~\ref{fig:abl-count-N}.}
\label{fig:abl-zi-ATE}
\end{figure}

\subsection{Cross-cutting interpretation}\label{sec:emp-interp}

Three threads run through the reference-cell results and the three
ablation sweeps.

\paragraph{Treating counts as Gaussian systematically under-states
posterior uncertainty.} Across all $56$ count-only and $56$
zero-inflated cells (excluding the single missing cell), Gaussian
BCF intervals are $25$--$60\%$ narrower than CountBCF's at the same
cell, yet their empirical coverage sits $10$--$15$ percentage
points below the nominal $0.95$ on most cells and below $0.55$ in
the worst (linear ZINB at $p = 250$, ATE coverage $0.34$). This is
the worst kind of miscalibration for a policy report: a $0.95$
label on an interval that empirically covers $0.55$ implies the
analyst is communicating an order-of-magnitude over-confident
posterior. CountBCF's wider intervals are not a free-lunch loss in
sharpness --- they are intervals that actually cover, and in the
zero-inflated case they cover at essentially the nominal level.

\paragraph{The CountBCF / Gaussian gap widens with likelihood
mis-specification.} The relative PEHE improvement of CountBCF over
BCF (Gaussian) is monotone in the severity of the count likelihood
mis-specification: smallest on linear Poisson (the cell closest to
Gaussian) at $-23\%$, larger on the nonlinear and / or NB cells,
and largest on the nonlinear ZINB cell at $-54\%$. Comparable
ordering holds on RMSE(ATE) and on the gap between empirical and
nominal coverage. This is direct evidence that the source of the
improvement is the count likelihood --- the only structural
difference between the two methods --- rather than incidental
hyperparameter or post-processing choices.

\paragraph{The two methods react quite differently to nuisance
covariates.} CountBCF's PEHE and coverage are nearly flat across
$p \in \{5, 50, 250\}$ at $N = 250$ on both sub-studies, while the
Gaussian BCF degrades steadily. We attribute this to the
log-scale regularisation: on the log scale, irrelevant covariates
have a hard time standing in for residual count heteroscedasticity
because the Dirichlet-style splitting rules concentrate on the few
covariates that improve the count log-likelihood, and the
Metropolis acceptance ratio under the count likelihood penalises
spurious splits more aggressively than the same prior under a
homoscedastic-Gaussian likelihood of the same nominal mean. The
gap is a regularisation gap, not a propensity-adjustment gap:
both methods use the propensity-as-covariate fix, and both share
the BCF $(\mu, \tau)$ decomposition. CountBCF additionally
benefits, on the ZI cells, from the channel decomposition: the
treatment effect on the structural-zero probability has its own
forest pair, so heterogeneity that the Gaussian baseline must
absorb into a single response-scale CATE is allocated to its
proper channel.

\subsection{Caveats}\label{sec:emp-caveats}

\begin{itemize}
\item \emph{Monte Carlo size.} With $N_{\mathrm{sim}} = 100$ per
  cell, cross-replicate SDs are the right scale for the variance
  bars in Figures~\ref{fig:abl-count-N}--\ref{fig:abl-zi-ATE}, but
  cell-level means still carry Monte Carlo noise of order
  $\widehat\sigma / \sqrt{100}$. Tight comparisons (e.g.\ the
  $1.21$ vs.\ $1.25$ RMSE on \verb|linear_nb_cate|) are within
  that noise.
\item \emph{ATE coverage is replicate-realised.} Coverage of the
  ATE is read off the \emph{realised} ATE for that replicate, not
  a single population value. This is the right way to score a
  posterior interval against the truth that the data actually
  speak to, but it does mean that ATE coverage moves with the
  realised variability of $\tau_0$ across the calibration target.
\item \emph{Missing count cell.} CountBCF on
  \verb|nonlinear_poisson_cate| at $N = 250$, $p = 250$ did not
  complete (the runner was killed by Colab's memory ceiling on
  that single cell). The corresponding curve in
  Figure~\ref{fig:abl-count-P} therefore shows only three of the
  four DGPs at $p = 250$ for CountBCF; the BCF (Gaussian) curve
  at the same cell completed normally. The aggregate $-34\%$ and
  $+29\%$ numbers for CountBCF in the text are averaged over the
  three completed DGPs at $p = 250$.
\item \emph{Large-$p$ PEHE variability.} The cross-replicate PEHE
  SD at $p = 250$ remains large for both methods on both
  sub-studies; cell-to-cell noise in the $p$-ablation panels at
  the high end should be interpreted with a Monte Carlo error of
  order $\widehat\sigma_{\mathrm{PEHE}} / \sqrt{100}$ in mind.
\item \emph{True propensity used.} All runs use
  $\widehat{\pi} = \pi$, the true propensity. This benchmark
  isolates the outcome model; real-data performance with an
  estimated $\widehat{\pi}$ will be worse for both methods to a
  similar degree, because the propensity-as-covariate adjustment
  is shared.
\item \emph{Two-method comparison.} CountBART, the S-learner
  built on the same log-linear backbone, was originally part of
  the benchmark but is dropped here on the memory grounds noted
  in Section~\ref{sec:method-methods}; an external comparison
  against generalised random forests, $X$-learner, and DR-learner
  on the same DGPs is left to future work.
\end{itemize}

% =======================================================================
\section{Discussion}\label{sec:discussion}
\noindent\emph{[Reserved.  To be written after the empirical
evaluations.  Suggested topics: applied dataset picks (count cells --
hospital admissions, citation counts; ZI cells -- healthcare
utilisation under a wellness intervention, insurance claims under
policy change, ecological abundance under habitat intervention,
scRNAseq under a perturbation), identifiability remarks
(\citealt{murray2021loglinear} \S4.2.2 for prior-induced
identification of $\fzeroSym, \foneSym$;
\citealt{li2012zipidentifiability} on ZIP identifiability conditions;
\citealt{wager2018estimation} on count-forest asymptotics),
reproducibility statement.]}

% =======================================================================
\bibliographystyle{plainnat}
\bibliography{refs}

\end{document}
